\documentclass[12pt]{article}
\usepackage[ukrainian,shorthands=off]{babel}   % shorthands=off: символ " не активний (потрібно для міток "..." у TikZ всередині \zadtask)
\usepackage{fontspec}
% --- НАЛАШТУВАННЯ ШРИФТІВ ---
% Шрифти Schoolbook-*.otf лежать у корені проєкту. Шлях підбирається автоматично:
% ../ (компіляція з папки теми), ./ (корінь проєкту), ../../ (підпапки теми 28); інакше -- TeX Gyre Schola.
\newcommand{\nmtSchoolbook}[1]{\setmainfont{Schoolbook-Regular.otf}[
    Path = #1,
    BoldFont = Schoolbook-Bold.otf,
    ItalicFont = Schoolbook-Italic.otf,
    BoldItalicFont = Schoolbook-BoldItalic.otf
]}
\IfFileExists{../Schoolbook-Regular.otf}{\nmtSchoolbook{../}}{%
 \IfFileExists{./Schoolbook-Regular.otf}{\nmtSchoolbook{./}}{%
  \IfFileExists{../../Schoolbook-Regular.otf}{\nmtSchoolbook{../../}}{%
   \setmainfont{texgyreschola}[Extension=.otf, UprightFont=*-regular, BoldFont=*-bold, ItalicFont=*-italic, BoldItalicFont=*-bolditalic]}}}
\usepackage[italic]{mathastext}
\usepackage[a4paper,margin=1.8cm,bottom=2cm,top=2cm]{geometry}
\usepackage{amsmath,amssymb}
\usepackage{tikz}
\usepackage{pgfplots}
\pgfplotsset{compat=1.18}
\usepgfplotslibrary{fillbetween}
\usetikzlibrary{calc,patterns,angles,quotes,intersections,babel,3d,shapes.symbols,shapes.geometric,decorations.pathreplacing,shadings,arrows.meta,hobby}
% ===================== ДИЗАЙН-СИСТЕМА РИСУНКІВ (scripts/ПРАВИЛА_ОФОРМЛЕННЯ.md, scripts/дизайн_завдань.py) =====================
\definecolor{nmtInk}{HTML}{1A1A1A}
\definecolor{nmtGrid}{HTML}{3A3A3A}      % сітка, як у НМТ 2026: темна й тонка (0,35 pt), клітинка 0,5 см
\definecolor{nmtBlue}{HTML}{1F6FB2}\definecolor{nmtGreen}{HTML}{2E8B57}\definecolor{nmtRed}{HTML}{C0392B}
\definecolor{nmtOrange}{HTML}{D9822B}\definecolor{nmtViolet}{HTML}{7D4CA8}
\definecolor{nmtFillA}{HTML}{D3E5F5}\definecolor{nmtFillB}{HTML}{D8EDD5}\definecolor{nmtFillC}{HTML}{FBE1C8}
\definecolor{nmtFillD}{HTML}{E4E4E4}\definecolor{nmtFillE}{HTML}{E7DCF2}
\tikzset{
  ultra thin/.style={line width=0.2pt}, very thin/.style={line width=0.3pt}, thin/.style={line width=0.45pt},
  semithick/.style={line width=0.6pt}, thick/.style={line width=0.75pt}, very thick/.style={line width=1.1pt},
  ultra thick/.style={line width=1.5pt},
  >={Stealth[length=2.3mm, width=1.7mm]},
  every picture/.append style={line join=round, line cap=round},
  nmt dot/.style={circle, fill=nmtInk, draw=none, inner sep=0pt, minimum size=3.6pt},
  nmt open dot/.style={circle, fill=white, draw=nmtInk, line width=0.6pt, inner sep=0pt, minimum size=4.2pt},
  nmt grid/.style={draw=nmtGrid, line width=0.35pt},
  nmt curve/.style={line width=1.1pt},
}
% вирази: ≤ і ≥ -- похилі, як у НМТ (десяткову кому без пробілу вже дає mathastext)
\let\leq\leqslant \let\le\leqslant \let\geq\geqslant \let\ge\geqslant
\pgfplotsset{every axis/.append style={axis line style={line width=0.6pt}, tick style={line width=0.45pt},
  tick label style={font=\footnotesize}, label style={font=\small}, grid style={nmtGrid, line width=0.35pt}}}
% ===================== кінець дизайн-системи =====================
% підпис кута біля дуги: \nmtAngleLabel{A}{B}{C}{радіус дуги, см}{підпис} -- кут AB C проти годинникової стрілки, як \pic{angle=A--B--C}.
% Підпис ставиться на бісектрисі, найближчим краєм на 0,08 см від дуги (радіус дуги масштабується разом з рисунком, підпис -- ні),
% тож не відлітає від кута, хоч який масштаб рисунка й хоч яка довжина підпису
\newcommand{\nmtAngleLabel}[5]{%
  \path let \p1=(#2), \p2=(#1), \p3=(#3), \p4=(0,0), \p5=(1,0),
    \n1={atan2(\y2-\y1,\x2-\x1)}, \n2={atan2(\y3-\y1,\x3-\x1)},
    \n3={mod(\n2-\n1+720,360)}, \n4={\n1+\n3/2}, \n5={\n4+180},
    \n6={veclen(\x5-\x4,\y5-\y4)/28.45274}, \n7={(#4*\n6+0.08)*28.45274}
  in (#2) node[anchor=\n5, inner sep=1.2pt, shift={(\n4:\n7)}] {#5};}

\usepackage{xcolor,array,fancyhdr,enumitem,multicol}
\usepackage{colortbl,multirow,diagbox,needspace}
\usepackage{tcolorbox}
\tcbuselibrary{skins,breakable}
\usepackage[hidelinks]{hyperref}
% водяний знак; щоб зібрати без нього (версія для вчителів), перед \documentclass пишуть \def\nmtnowatermark{}
\ifdefined\nmtnowatermark\else
\usepackage[firstpage=false, text=@pvtr2525, scale=0.6, color=gray!12]{draftwatermark}
\fi

% --- АВТОСТИСКАННЯ: рисунок/сітка, ширша за колонку, масштабується до ширини колонки ---
\newsavebox{\nmtfitbox}
\newenvironment{nmtfit}{\begin{lrbox}{\nmtfitbox}}{\end{lrbox}%
  \ifdim\wd\nmtfitbox>\linewidth \resizebox{\linewidth}{!}{\usebox{\nmtfitbox}}\else\usebox{\nmtfitbox}\fi}
\let\nmtIncludegraphicsOrig\includegraphics
\renewcommand{\includegraphics}[2][]{\begin{nmtfit}\nmtIncludegraphicsOrig[#1]{#2}\end{nmtfit}}


\definecolor{mainGreen}{RGB}{34, 120, 64}
\definecolor{yearOrange}{RGB}{220, 100, 30}
\definecolor{yearcolor}{RGB}{220, 100, 30}
\definecolor{titleBg}{RGB}{225, 240, 220}
\definecolor{instrBg}{RGB}{255, 240, 220}
\definecolor{instrBorder}{RGB}{220, 150, 80}
\definecolor{headerblue}{RGB}{0, 102, 204}   % використовується в рисунках старої бази

\setlength{\parindent}{0pt}
\emergencystretch=1.5em

\pagestyle{fancy}
\fancyhf{}
\renewcommand{\headrulewidth}{0.4pt}
\renewcommand{\footrulewidth}{0.4pt}
\fancyhead[C]{\small\color{gray!80} Вирази: завдання на відповідність}
\fancyhead[R]{\small\color{mainGreen}\textbf{\href{https://t.me/pvtr2525}{@pvtr2525}}}
\fancyfoot[L]{\small\color{gray!80}\href{https://t.me/pvtr2525}{ТГ: @pvtr2525}}
\fancyfoot[C]{\small\color{gray!80}\thepage}
\fancyfoot[R]{\small\color{gray!80}НМТ 2022--2026}

% --- ТАБЛИЦІ ВІДПОВІДЕЙ ---
% ширина клітинки = (ширина рядка - відступи) / 5: на всю сторінку це ~3 см, у minipage поруч із рисунком таблиця стискається;
% п'ять колонок |m{w}| мають 10 відступів \tabcolsep і 6 вертикальних ліній -- тоді таблиця рівно на \linewidth
\newlength{\nmtcellw}
\newcommand{\nmtSetCell}{\setlength{\nmtcellw}{\dimexpr(\linewidth-10\tabcolsep-6\arrayrulewidth)/5\relax}}
\newcommand{\answerTable}[5]{%
\par\nopagebreak\vspace{0.15cm}
\noindent\nmtSetCell
\begin{tabular}{|*{5}{>{\centering\arraybackslash}m{\nmtcellw}|}}
\hline
\rule[-0.15cm]{0pt}{0.6cm}\textbf{А} & \rule[-0.15cm]{0pt}{0.6cm}\textbf{Б} & \rule[-0.15cm]{0pt}{0.6cm}\textbf{В} & \rule[-0.15cm]{0pt}{0.6cm}\textbf{Г} & \rule[-0.15cm]{0pt}{0.6cm}\textbf{Д} \\
\hline
\rule[-0.2cm]{0pt}{0.75cm}#1 & \rule[-0.2cm]{0pt}{0.75cm}#2 & \rule[-0.2cm]{0pt}{0.75cm}#3 & \rule[-0.2cm]{0pt}{0.75cm}#4 & \rule[-0.2cm]{0pt}{0.75cm}#5 \\
\hline
\end{tabular}
\par\vspace{0.25cm}
}
\newcommand{\answerTableTall}[5]{%
\par\nopagebreak\vspace{0.15cm}
\noindent\nmtSetCell
\begin{tabular}{|*{5}{>{\centering\arraybackslash}m{\nmtcellw}|}}
\hline
\rule[-0.2cm]{0pt}{0.7cm}\textbf{А} & \rule[-0.2cm]{0pt}{0.7cm}\textbf{Б} & \rule[-0.2cm]{0pt}{0.7cm}\textbf{В} & \rule[-0.2cm]{0pt}{0.7cm}\textbf{Г} & \rule[-0.2cm]{0pt}{0.7cm}\textbf{Д} \\
\hline
\rule[-0.45cm]{0pt}{1.2cm}#1 & \rule[-0.45cm]{0pt}{1.2cm}#2 & \rule[-0.45cm]{0pt}{1.2cm}#3 & \rule[-0.45cm]{0pt}{1.2cm}#4 & \rule[-0.45cm]{0pt}{1.2cm}#5 \\
\hline
\end{tabular}
\par\vspace{0.25cm}
}
\newcommand{\instructionBox}[1]{%
\par\vspace{0.3cm}
\begin{tcolorbox}[colback=instrBg, colframe=instrBorder, boxrule=0.6pt, arc=2pt,
    left=8pt, right=8pt, top=4pt, bottom=4pt]
\centering\small\bfseries #1
\end{tcolorbox}
\par\vspace{0.15cm}
}
\newcommand{\sectionTitle}[1]{%
\par\vspace{0.4cm}
\begin{tcolorbox}[colback=titleBg, colframe=titleBg, boxrule=0pt, arc=8pt,
    halign=center, fontupper=\Large\bfseries\color{mainGreen}]
\hyphenpenalty=10000 \exhyphenpenalty=10000 #1
\end{tcolorbox}
\par\vspace{0.3cm}
}
\newcommand{\task}[2]{\noindent\textbf{#1.}\ \ #2\par\vspace{0.2cm}}
% --- НМТ-СТИЛЬ ПОЛЕ ВІДПОВІДІ ---
\newcommand{\nmtAnswerBox}{%
\par\nopagebreak\vspace{0.2cm}\noindent
Відповідь:\ \
\begin{tabular}{@{}*{4}{|>{\centering\arraybackslash}p{0.8cm}}|@{\hspace{0.1cm}\textbf{,}\hspace{0.1cm}}*{3}{|>{\centering\arraybackslash}p{0.8cm}}|@{}}
\hline
\rule[-0.2cm]{0pt}{0.7cm} & & & & & & \\
\hline
\end{tabular}
\par\vspace{0.3cm}
}
% --- СІТКА ДЛЯ МАТЧИНГУ ---
\newsavebox{\nmtgridbox}
\newcommand{\matchingGrid}{%
    \sbox{\nmtgridbox}{\matchingGridRaw}%
    \ifdim\wd\nmtgridbox>\linewidth \resizebox{\linewidth}{!}{\usebox{\nmtgridbox}}\else\usebox{\nmtgridbox}\fi}
\newcommand{\matchingGridRaw}{%
    \begingroup
    \renewcommand{\arraystretch}{1.15}
    \setlength{\tabcolsep}{4pt}
    \footnotesize
    \begin{tabular}{r|c|c|c|c|c|}
         \multicolumn{1}{c}{} & \multicolumn{1}{c}{\textbf{А}} & \multicolumn{1}{c}{\textbf{Б}} & \multicolumn{1}{c}{\textbf{В}} & \multicolumn{1}{c}{\textbf{Г}} & \multicolumn{1}{c}{\textbf{Д}} \\ \cline{2-6}
         \textbf{1} & \phantom{X} & \phantom{X} & \phantom{X} & \phantom{X} & \phantom{X} \\ \cline{2-6}
         \textbf{2} & & & & & \\ \cline{2-6}
         \textbf{3} & & & & & \\ \cline{2-6}
    \end{tabular}
    \endgroup
}
% --- ДОДАТКОВІ МАКРОСИ ЗБІРНИКА ---
\newcommand{\taskBlock}[1]{%
\par\vspace{0.45cm}
\noindent{\color{mainGreen}\rule{\linewidth}{0.8pt}}\par\vspace{0.12cm}
\noindent{\small\bfseries\color{mainGreen}#1}\par\vspace{0.25cm}
}
\newcounter{zad}
\newcommand{\zadtask}[1]{\stepcounter{zad}\noindent\textbf{\thezad.}\ \ #1\par\nopagebreak\vspace{0.2cm}}
\newcommand{\zadnum}{\stepcounter{zad}\textbf{\thezad.}\ \ }
\newcounter{ans}
\newcommand{\ansitem}[1]{\stepcounter{ans}\noindent\textbf{\theans.}\ #1\par\vspace{0.07cm}}
% мітка року: нерозривна, притиснута праворуч; якщо не вміщається -- праворуч на новому рядку
\newcommand{\nmtyear}[1]{\unskip\nobreak\finalhyphendemerits=0 \hfill\penalty100\hskip1em\hbox{}\nobreak\hfill{\small\color{yearOrange}\mbox{\rule{0pt}{2.3ex}(НМТ~#1)}}}
% --- МАКРОС КОРОТКОГО РОЗВ'ЯЗКУ ---
\newcommand{\solution}[1]{%
\par\vspace{0.12cm}
\begin{tcolorbox}[colback=titleBg!45, colframe=mainGreen!55, boxrule=0.4pt, arc=2pt,
    left=6pt, right=6pt, top=3pt, bottom=3pt, breakable]
\footnotesize\textbf{Короткий розв'язок.}\ #1
\end{tcolorbox}
\par\vspace{0.12cm}
}
% Розв'язки й відповіді (є до завдань НМТ-2026): \showsolutionstrue -- показати, \showsolutionsfalse -- сховати
\newif\ifshowsolutions
\showsolutionsfalse
% --- ЗАГОЛОВКИ З ПУНКТАМИ КЛІКАБЕЛЬНОГО ЗМІСТУ ---
\setcounter{tocdepth}{2}
\newcommand{\chapterTitle}[1]{%
\clearpage\phantomsection\addcontentsline{toc}{section}{#1}%
\sectionTitle{#1}}
\newcommand{\typeTitle}[1]{%
\needspace{6\baselineskip}%
\phantomsection\addcontentsline{toc}{subsection}{\quad #1}%
\par\vspace{0.3cm}
\noindent{\large\bfseries\color{yearOrange}#1}\par\vspace{0.08cm}
\noindent{\color{yearOrange}\rule{\linewidth}{0.4pt}}\par\nopagebreak\vspace{0.2cm}}
\raggedbottom
% усередині одного завдання сторінка не розривається: нерозривний вертикальний проміжок
% і нерозривна межа абзаців (умова ніколи не відділяється від варіантів, рисунка чи сітки)
\newcommand{\nbvspace}[1]{\nopagebreak\vspace{#1}}
\newcommand{\nmtnobreak}{\ifvmode\penalty10000 \fi}
% --- ЄДИНИЙ МАКЕТ ЗАВДАНЬ НА ВІДПОВІДНІСТЬ ---
% мітка в колонці сталої ширини, текст -- у parbox: продовження довгого рядка
% вирівнюється під текстом, а не під міткою; відступ між пунктами однаковий скрізь
\newlength{\nmtMatchLab}\setlength{\nmtMatchLab}{1.6em}
\newlength{\nmtMatchSep}\setlength{\nmtMatchSep}{0.28cm}
\newcommand{\matchHead}[1]{\par\noindent\textit{#1}\par\nopagebreak\vspace{0.2cm}}
% шаблонний заголовок стовпця зі старого макета -- лише коли завдання не має власного \matchHead
\makeatletter\newcommand{\nmtHeadUnless}[2]{\in@{\matchHead}{#1}\ifin@\else#2\fi}\makeatother
\newcommand{\matchItem}[2]{\par\noindent\makebox[\nmtMatchLab][l]{\textbf{#1}}%
\parbox[t]{\dimexpr\linewidth-\nmtMatchLab\relax}{\raggedright #2}\par\nopagebreak\vspace{\nmtMatchSep}}
\newcommand{\ansTheme}[1]{\par\vspace{0.25cm}\noindent{\bfseries\color{mainGreen}#1}\par\vspace{0.12cm}}
\newcommand{\ansType}[1]{\par\vspace{0.1cm}\noindent{\itshape\color{yearOrange}#1}\par\vspace{0.08cm}}

% варіанти відповіді окремими рядками (довгі вирази не влазять у клітинки таблиці)
\newcommand{\answerRows}[5]{\par\nopagebreak\vspace{0.15cm}%
\matchItem{А}{#1}\matchItem{Б}{#2}\matchItem{В}{#3}\matchItem{Г}{#4}\matchItem{Д}{#5}}
\graphicspath{{./}{1. Числа. Звичайні та десяткові дроби. Модуль числа/}{26. Тригонометричні вирази/}{28. Показникова функція. Показникові рівняння/Показникові функція і вирази/}{3. Степені. Одночлени. стандартний вигляд числа/}{30. Логарифм. Логарифмічна функція/}{4.Многочлени. ФСМ/}{5. дробові раціональні вирази і перетворення/}{7. Квадратні корені. Корені вищих степенів. Раціональний степінь/}}

\begin{document}
\begingroup\setcounter{zad}{0}
% --- макроси з файлів сесій НМТ-2026 (для рисунків) ---
\renewcommand{\headrulewidth}{0.4pt}
\renewcommand{\footrulewidth}{0.4pt}
\definecolor{gold}{RGB}{240, 190, 40}
\definecolor{silver}{RGB}{170, 170, 175}
\definecolor{bronze}{RGB}{190, 120, 70}
\definecolor{hexcol}{RGB}{200,110,60}
\newcommand{\hexthree}[2]{%
    \draw[hexcol, line width=1.1pt]
        ($(#1,#2)+(0:30)$) -- ($(#1,#2)+(60:30)$) -- ($(#1,#2)+(120:30)$) --
        ($(#1,#2)+(180:30)$) -- ($(#1,#2)+(240:30)$) -- ($(#1,#2)+(300:30)$) -- cycle;
}

% --- сумісність зі старою базою (діє, якщо тема не має власного визначення) ---
\providecommand{\trigStyle}[1]{#1}
\providecommand{\answerTableSmall}[5]{%
\begingroup\setlength{\tabcolsep}{2pt}\small\nmtSetCell
\begin{tabular}{|*{5}{>{\centering\arraybackslash}m{\nmtcellw}|}}
\hline
\rule[-0.2cm]{0pt}{0.6cm}\textbf{А} & \textbf{Б} & \textbf{В} & \textbf{Г} & \textbf{Д} \\
\hline
\rule[-0.4cm]{0pt}{0.9cm}#1 & \rule[-0.4cm]{0pt}{0.9cm}#2 & \rule[-0.4cm]{0pt}{0.9cm}#3 & \rule[-0.4cm]{0pt}{0.9cm}#4 & \rule[-0.4cm]{0pt}{0.9cm}#5 \\
\hline
\end{tabular}\endgroup}
\providecommand{\matchingLayout}[3]{%
    \noindent
    \begin{minipage}[t]{0.36\textwidth}\vspace{0pt}\raggedright
        #1
    \end{minipage}%
    \hfill
    \begin{minipage}[t]{0.43\textwidth}\vspace{0pt}\raggedright
        #2
    \end{minipage}%
    \hfill
    \begin{minipage}[t]{0.19\textwidth}
        \vspace{0pt}
        \begin{flushright}
        #3
        \end{flushright}
    \end{minipage}}
\providecommand{\matchingLayoutBottom}[2]{%
    \noindent
    \begin{minipage}[t]{0.48\textwidth}\vspace{0pt}\raggedright
        #1
    \end{minipage}%
    \hfill
    \begin{minipage}[t]{0.48\textwidth}\vspace{0pt}\raggedright
        #2
    \end{minipage}
    \vspace{0.5cm}
    \noindent
    \matchingGrid}
\providecommand{\answerListVertical}[5]{%
    \par\nopagebreak\vspace{0.2cm}
    \begin{itemize}[itemsep=0.4cm, leftmargin=1.5cm, labelsep=0.5cm]
        \item[\textbf{А}] #1
        \item[\textbf{Б}] #2
        \item[\textbf{В}] #3
        \item[\textbf{Г}] #4
        \item[\textbf{Д}] #5
    \end{itemize}
    \vspace{0.2cm}}

\graphicspath{{./}{11. Трикутники/}}
\setcounter{zad}{0}

\chapterTitle{Вирази: завдання на відповідність}

\noindent{\small Збірник містить \textbf{усі 73 завдань} НМТ 2022--2026 років на встановлення відповідності з бази (вирази). \textbf{Частина~1} (5) -- без тригонометрії, показникових виразів і логарифмів, \textbf{частина~2} (68) -- решта. Усередині частин завдання згруповано за типом; біля кожного вказано рік. Відповіді~--- на останній сторінці.\par\vspace{0.15cm}
За роками: \mbox{2022~--- 1}, \mbox{2023~--- 20}, \mbox{2024~--- 18}, \mbox{2025~--- 17}, \mbox{2026~--- 17}.}\par\vspace{0.45cm}

\begin{samepage}
\needspace{0.45\textheight}\sectionTitle{Частина 1. Без тригонометрії, показникових виразів і логарифмів}
\typeTitle{Значення числового виразу}
% НМТ 2026, сесія 11.06, завдання №16 --- основна тема: Квадратні корені. Корені вищих степенів. Раціональний степінь
\zadtask{Узгодьте вираз \mbox{(1--3)} із його значенням \mbox{(А--Д)}. \nmtyear{2026}\par
\nopagebreak\nbvspace{0.3cm}
\noindent
\matchingLayout{
\matchHead{Вираз}
\matchItem{1}{$\left|\dfrac{2}{5} - \dfrac{12}{3}\right|$}
\matchItem{2}{$\dfrac{\sqrt{24}}{\sqrt{6}}$}
\matchItem{3}{$5^3 \cdot 0{,}2^4$}
}{
\matchHead{Значення виразу}
\matchItem{А}{$0{,}2$}
\matchItem{Б}{$2$}
\matchItem{В}{$3{,}6$}
\matchItem{Г}{$4$}
\matchItem{Д}{$5$}
}{
\matchingGrid
}
}
\par\penalty-20
\end{samepage}

\begin{samepage}
\needspace{0.45\textheight}\typeTitle{Проміжок, якому належить значення числового виразу}
% Завдання 8 (на відповідність)
\zadtask{Установіть відповідність між виразом \mbox{(1--3)} та проміжком \mbox{(А--Д)}, якому належить значення цього виразу, якщо $a = -0{,}5$. \nmtyear{2023}}

\nmtnobreak
\nopagebreak\nbvspace{0.3cm}
\noindent
\matchingLayout{
\matchHead{Вираз}
\matchItem{1}{$|a|$}
\matchItem{2}{$a^3$}
\matchItem{3}{$\dfrac{1}{a}$}
}{
\matchHead{Проміжок}
\matchItem{А}{$(-\infty; -2)$}
\matchItem{Б}{$[-2; -1)$}
\matchItem{В}{$[-1; 0)$}
\matchItem{Г}{$[0; 1)$}
\matchItem{Д}{$[1; +\infty)$}
}{
\matchingGrid
}
\par\penalty-20\vspace{0.25cm}
\end{samepage}

\begin{samepage}
\needspace{0.45\textheight}\typeTitle{Твердження про значення виразу:\\ натуральне, ціле, раціональне, ірраціональне}
% Завдання 20 (на відповідність)
\zadtask{Установіть відповідність між виразом \mbox{(1--3)} та твердженням про його значення \mbox{(А--Д)}, яке є правильним. \nmtyear{2024}}

\nmtnobreak
\nopagebreak\nbvspace{0.3cm}
\noindent
\matchingLayout{
\matchHead{Вираз}
\matchItem{1}{$(\sqrt{3} - 1)^2$}
\matchItem{2}{$\sqrt[3]{-8^2}$}
\matchItem{3}{$\dfrac{\sqrt{12}}{\sqrt{3}}$}
}{
\matchHead{Твердження про значення виразу}
\matchItem{А}{є ірраціональним додатним числом}
\matchItem{Б}{є ірраціональним від'ємним числом}
\matchItem{В}{дорівнює 0}
\matchItem{Г}{є натуральним числом}
\matchItem{Д}{є цілим від'ємним числом}
}{
\matchingGrid
}
\par\penalty-20\vspace{0.25cm}
\end{samepage}

\begin{samepage}
% Завдання 39 (на відповідність)
\zadtask{Узгодьте вираз \mbox{(1--3)} із твердженням про його значення \mbox{(А--Д)}, що є правильним для цього виразу. \nmtyear{2025}}

\nmtnobreak
\nopagebreak\nbvspace{0.3cm}
\matchingLayout{
\matchHead{Вираз}
\matchItem{1}{$\sqrt[3]{-27}$}
\matchItem{2}{$\dfrac{\sqrt{12}}{\sqrt{300}}$}

\nmtnobreak
\matchItem{3}{$(\sqrt{3} - 1)^2 + 2\sqrt{3}$}
}{
\matchHead{Твердження про значення виразу}
\matchItem{А}{є цілим від'ємним числом}
\matchItem{Б}{є ірраціональним додатним числом}
\matchItem{В}{є ірраціональним від'ємним числом}
\matchItem{Г}{є натуральним числом}
\matchItem{Д}{є раціональним нецілим числом}
}{
\matchingGrid
}
\par\penalty-20\vspace{0.25cm}
\end{samepage}

\begin{samepage}
\needspace{0.45\textheight}\typeTitle{Тотожно рівний вираз зі змінною}
% НМТ 2026, сесія 15.06, завдання №16 --- основна тема: Квадратні корені. Корені вищих степенів. Раціональний степінь
\zadtask{Узгодьте вираз \mbox{(1--3)} з тотожно рівним йому виразом \mbox{(А--Д)}, якщо $a \neq 0$, $a \neq 1$. \nmtyear{2026}\par
\nopagebreak\nbvspace{0.3cm}
\noindent
\matchingLayout{
\matchHead{Вираз}
\matchItem{1}{$\sqrt{a^2}$}
\matchItem{2}{$a^3 : a^4$}
\matchItem{3}{$\dfrac{a^2 - a}{1 - a}$}
}{
\matchHead{Тотожно рівний вираз}
\matchItem{А}{$a$}
\matchItem{Б}{$2a$}
\matchItem{В}{$\dfrac{1}{a}$}
\matchItem{Г}{$-a$}
\matchItem{Д}{$|a|$}
}{
\matchingGrid
}
}
% Відповідь: 1--Д, 2--В, 3--Г
\par\penalty-20
\end{samepage}

\begin{samepage}
\needspace{0.45\textheight}\sectionTitle{Частина 2. З тригонометрією, показниковими виразами чи логарифмами}
\typeTitle{Значення числового виразу}
% Завдання 21 (на відповідність)
\zadtask{Установіть відповідність між виразом \mbox{(1--3)} та значенням \mbox{(А--Д)} цього виразу. \nmtyear{2024}}

\nmtnobreak
\nopagebreak\nbvspace{0.3cm}
\matchingLayout{
\matchHead{Вираз}
\matchItem{1}{$\dfrac{3^{-5}}{3^{-6}}$}

\nmtnobreak
\matchItem{2}{$\log_2 0{,}1 + \log_2 320$}
\matchItem{3}{$4\cos^2 30° - 4\sin^2 30°$}
}{
\matchHead{Значення виразу}
\matchItem{А}{$1$}
\matchItem{Б}{$2$}
\matchItem{В}{$3$}
\matchItem{Г}{$4$}
\matchItem{Д}{$5$}
}{
\matchingGrid
}
\par\penalty-20\vspace{0.25cm}
\end{samepage}

\begin{samepage}
% Завдання 28 (на відповідність)
\zadtask{Доберіть до числового виразу \mbox{(1--3)} його значення \mbox{(А--Д)}. \nmtyear{2025}}

\nmtnobreak
\nopagebreak\nbvspace{0.3cm}
\noindent
\matchingLayout{
\matchHead{Вираз}
\matchItem{1}{$\log_3 27$}
\matchItem{2}{$\mathrm{tg}\,\dfrac{2\pi}{3}$}
\matchItem{3}{$\dfrac{1}{\sqrt{5} - \sqrt{2}} : (\sqrt{5} + \sqrt{2})$}
}{
\matchHead{Значення виразу}
\matchItem{А}{$-3$}
\matchItem{Б}{$\sqrt{3}$}
\matchItem{В}{$\dfrac{1}{3}$}
\matchItem{Г}{$-\sqrt{3}$}
\matchItem{Д}{$3$}
}{
\matchingGrid
}
\par\penalty-20\vspace{0.25cm}
\end{samepage}

\begin{samepage}
% Завдання 40 (на відповідність)
\zadtask{Узгодьте вираз \mbox{(1--3)} із його значенням \mbox{(А--Д)}. \nmtyear{2025}}

\nmtnobreak
\nopagebreak\nbvspace{0.3cm}
\matchingLayout{
\matchHead{Вираз}
\matchItem{1}{$\sin\left(-\dfrac{\pi}{6}\right)$}
\matchItem{2}{$\pi - |\pi + 2|$}
\matchItem{3}{$\left(-\dfrac{1}{\sqrt{2}}\right)^2$}
}{
\matchHead{Значення виразу}
\matchItem{А}{$-\dfrac{1}{2}$}
\matchItem{Б}{$\dfrac{1}{2}$}
\matchItem{В}{$2$}
\matchItem{Г}{$-2$}
\matchItem{Д}{$\sqrt{2}$}
}{
\matchingGrid
}
\par\penalty-20\vspace{0.25cm}
\end{samepage}

\begin{samepage}
% Завдання 35 (на відповідність)
\zadtask{Доберіть до числового виразу \mbox{(1--3)} його значення \mbox{(А--Д)}. \nmtyear{2025}}

\nmtnobreak
\nopagebreak\nbvspace{0.3cm}
\noindent
\matchingLayout{
\matchHead{Вираз}
\matchItem{1}{$\lg 100$}
\matchItem{2}{$\sin\dfrac{3\pi}{2}$}
\matchItem{3}{$(2 + \sqrt{3})(2 - \sqrt{3})$}
}{
\matchHead{Значення виразу}
\matchItem{А}{$-1$}
\matchItem{Б}{$0$}
\matchItem{В}{$1$}
\matchItem{Г}{$2$}
\matchItem{Д}{$10$}
}{
\matchingGrid
}
\par\penalty-20\vspace{0.25cm}
\end{samepage}

\begin{samepage}
% Завдання 3 (на відповідність)
\zadtask{Узгодьте вираз \mbox{(1--3)} із його значенням \mbox{(А--Д)}. \nmtyear{2025}}

\nmtnobreak
\nopagebreak\nbvspace{0.3cm}
\matchingLayout{
\matchHead{Вираз}
\matchItem{1}{$\dfrac{3}{3^{-3}}$}

\nmtnobreak
\matchItem{2}{$\log_8 \sqrt[3]{2}$}
\matchItem{3}{$2(\cos 30° - 0{,}5)(\cos 30° + 0{,}5)$}
}{
\matchHead{Значення виразу}
\matchItem{А}{$1$}
\matchItem{Б}{$\dfrac{1}{9}$}
\matchItem{В}{$\sqrt{3}$}
\matchItem{Г}{$3$}
\matchItem{Д}{$81$}
}{
\matchingGrid
}
\par\penalty-20\vspace{0.25cm}
\end{samepage}

\begin{samepage}
% НМТ 2026, сесія 17.06, завдання №16 --- також у темі: Квадратні корені. Корені вищих степенів. Раціональний степінь, Логарифмічні вирази
\zadtask{Узгодьте вираз \mbox{(1--3)} з його значенням \mbox{(А--Д)}. \nmtyear{2026}\par
\nopagebreak\nbvspace{0.3cm}
\noindent
\matchingLayout{
\matchHead{Вираз}
\matchItem{1}{$|-2| \cdot 2^0$}
\matchItem{2}{$\sqrt[3]{-\dfrac{1}{8}}$}
\matchItem{3}{$3^{-\log_3 2}$}
}{
\matchHead{Значення виразу}
\matchItem{А}{$-2$}
\matchItem{Б}{$-0{,}5$}
\matchItem{В}{$0$}
\matchItem{Г}{$0{,}5$}
\matchItem{Д}{$2$}
}{
\matchingGrid
}
}
% Відповідь: 1--Д, 2--Б, 3--Г
\par\penalty-20
\end{samepage}

\begin{samepage}
% НМТ 2026, сесія 5.06, завдання №16 --- також у темі: Логарифмічні вирази
\zadtask{Узгодьте вираз \mbox{(1--3)} із його значенням \mbox{(А--Д)}. \nmtyear{2026}

\nmtnobreak
\nopagebreak\nbvspace{0.3cm}
\noindent
\matchingLayout{
\matchHead{Вираз}
\matchItem{1}{$\left|\,-6 + \cos 4\pi\,\right|$}
\matchItem{2}{$\sqrt{2}\cdot\sqrt{18}$}
\matchItem{3}{$\log_{1/2} 32$}
}{
\matchHead{Значення}
\matchItem{А}{$-6$}
\matchItem{Б}{$-5$}
\matchItem{В}{$5$}
\matchItem{Г}{$6$}
\matchItem{Д}{$7$}
}{
\matchingGrid
}
}
% Відповідь: 1~--~В, 2~--~Г, 3~--~Б
\ifshowsolutions\solution{$1$: $|-6+\cos4\pi|=|-6+1|=5$ (В); $2$: $\sqrt{2}\cdot\sqrt{18}=\sqrt{36}=6$ (Г); $3$: $\log_{1/2}32=\log_{1/2}2^5=-5$ (Б). Відповідь: \textbf{1~--~В, 2~--~Г, 3~--~Б}.}\fi
\par\penalty-20
\end{samepage}

\begin{samepage}
% НМТ 2026, сесія 19.06, завдання №16 --- основна тема: Квадратні корені. Корені вищих степенів. Раціональний степінь
\zadtask{Доберіть до числового виразу \mbox{(1--3)} його значення \mbox{(А--Д)}. \nmtyear{2026}\par
\nopagebreak\nbvspace{0.3cm}
\noindent
\matchingLayout{
\matchHead{Числовий вираз}
\matchItem{1}{$\dfrac{8}{5} : \dfrac{4}{15}$}
\matchItem{2}{$\sqrt{2} \cdot \left(\sqrt{2}\right)^3$}
\matchItem{3}{$\log_{16} 2$}
}{
\matchHead{Значення виразу}
\matchItem{А}{$-\dfrac{1}{4}$}
\matchItem{Б}{$4$}
\matchItem{В}{$\dfrac{1}{4}$}
\matchItem{Г}{$6$}
\matchItem{Д}{$\dfrac{1}{6}$}
}{
\matchingGrid
}
}
% Відповідь: 1--Г,\ \ 2--Б,\ \ 3--В
\par\penalty-20
\end{samepage}

\begin{samepage}
% НМТ 2026, сесія 16.06, завдання №16 --- також у темі: Логарифмічні вирази
\zadtask{Установіть відповідність між виразом \mbox{(1--3)} та його значенням \mbox{(А--Д)}. \nmtyear{2026}\par
\nopagebreak\nbvspace{0.3cm}
\noindent
\matchingLayout{
\matchHead{Вираз}
\matchItem{1}{$\log_2 \dfrac{1}{16}$}
\matchItem{2}{$\dfrac{2}{\sqrt{0{,}25}}$}
\matchItem{3}{$(2^{-1})^2$}
}{
\matchHead{Значення виразу}
\matchItem{А}{$-4$}
\matchItem{Б}{$-\dfrac{1}{4}$}
\matchItem{В}{$\dfrac{1}{4}$}
\matchItem{Г}{$2$}
\matchItem{Д}{$4$}
}{
\matchingGrid
}
}
% Відповідь: 1--А; 2--Д; 3--В
\par\penalty-20
\end{samepage}

\begin{samepage}
% НМТ 2026, сесія 23.05, завдання №17
\zadtask{Установіть відповідність між виразом \mbox{(1--3)} та його значенням \mbox{(А--Д)}. \nmtyear{2026}

\nmtnobreak
\nopagebreak\nbvspace{0.3cm}

\nmtnobreak
\noindent
\matchingLayout{
\matchHead{Вираз}
\matchItem{1}{$\left|\dfrac{5}{3} - 2\right|$}
\matchItem{2}{$3^{-2} \cdot \left(\dfrac{1}{3}\right)^{-3}$}
\matchItem{3}{$\sqrt{3} \cdot \sin\dfrac{\pi}{3}$}
}{
\matchHead{Значення}
\matchItem{А}{$\dfrac{3}{2}$}
\matchItem{Б}{$3$}
\matchItem{В}{$\dfrac{1}{3}$}
\matchItem{Г}{$-\dfrac{1}{3}$}
\matchItem{Д}{$\dfrac{1}{2}$}
}{
\matchingGrid
}
}
% Відповідь: 1~--~В, 2~--~Б, 3~--~А
\ifshowsolutions\solution{$1)\ \left|\dfrac{5}{3}-2\right|=\left|-\dfrac{1}{3}\right|=\dfrac{1}{3}$~(В); $2)\ 3^{-2}\cdot 3^{3}=3^{1}=3$~(Б); $3)\ \sqrt{3}\cdot\sin\dfrac{\pi}{3}=\sqrt{3}\cdot\dfrac{\sqrt{3}}{2}=\dfrac{3}{2}$~(А). Відповідь: \textbf{1--В, 2--Б, 3--А}.}\fi
\par\penalty-20
\end{samepage}

\begin{samepage}
% НМТ 2026, сесія 10.06, завдання №16 --- також у темі: Логарифмічні вирази
\zadtask{Установіть відповідність між виразом \mbox{(1--3)}, де $\pi$~--- відома математична константа, та його значенням \mbox{(А--Д)}. \nmtyear{2026}\par
\nopagebreak\nbvspace{0.3cm}
\noindent
\matchingLayout{
\matchHead{Вираз}
\matchItem{1}{$\cos\left(-\dfrac{\pi}{3}\right)$}
\matchItem{2}{$\log_\pi \dfrac{1}{\pi}$}
\matchItem{3}{$|1 - \pi| - |\pi - 1|$}
}{
\matchHead{Значення виразу}
\matchItem{А}{$-1$}
\matchItem{Б}{$-0{,}5$}
\matchItem{В}{$0$}
\matchItem{Г}{$0{,}5$}
\matchItem{Д}{$2\pi$}
}{
\matchingGrid
}
}
\par\penalty-20
\end{samepage}

\begin{samepage}
\needspace{0.45\textheight}\typeTitle{Значення виразу за заданого значення змінної}
% Завдання 17 (на відповідність)
\zadtask{Установіть відповідність між виразом \mbox{(1--3)} та значенням \mbox{(А--Д)} цього виразу, якщо $x = \sqrt{5} - 4$. \nmtyear{2024}}

\nmtnobreak
\nopagebreak\nbvspace{0.3cm}
\noindent
\matchingLayout{
\matchHead{Вираз}
\matchItem{1}{$x^2 + 8x + 16$}
\matchItem{2}{$\dfrac{x - 1}{\sqrt{5}}$}

\nmtnobreak
\matchItem{3}{$\lg x^0$}
}{
\matchHead{Значення виразу}
\matchItem{А}{$5$}
\matchItem{Б}{$\sqrt{5}$}
\matchItem{В}{$0$}
\matchItem{Г}{$1 - \sqrt{5}$}
\matchItem{Д}{$-5$}
}{
\matchingGrid
}
\par\penalty-20\vspace{0.25cm}
\end{samepage}

\begin{samepage}
% Завдання 37 (на відповідність)
\zadtask{Доберіть до виразу \mbox{(1--3)} його значення \mbox{(А--Д)}, якщо $a = \sqrt{2}$. \nmtyear{2025}}

\nmtnobreak
\nopagebreak\nbvspace{0.3cm}
\noindent
\matchingLayout{
\matchHead{Вираз}
\matchItem{1}{$a^6$}
\matchItem{2}{$(a - 3)(a + 3)$}
\matchItem{3}{$81^{\log_3 a}$}
}{
\matchHead{Значення виразу}
\matchItem{А}{$-7$}
\matchItem{Б}{$-1$}
\matchItem{В}{$4$}
\matchItem{Г}{$8$}
\matchItem{Д}{$64$}
}{
\matchingGrid
}
\par\penalty-20\vspace{0.25cm}
\end{samepage}

\begin{samepage}
% НМТ 2026, сесія 18.06, завдання №16 --- основна тема: Квадратні корені. Корені вищих степенів. Раціональний степінь
\zadtask{Узгодьте вираз \mbox{(1--3)} із його значенням \mbox{(А--Д)}, якщо $a = \dfrac{3}{4}$. \nmtyear{2026}\par
\nopagebreak\nbvspace{0.3cm}
\noindent
\matchingLayout{
\matchHead{Вираз}
\matchItem{1}{$|a - 1| - a$}
\matchItem{2}{$\cos(2\pi a)$}
\matchItem{3}{$\sqrt{\dfrac{a}{3}}$}
}{
\matchHead{Значення виразу}
\matchItem{А}{$-\dfrac{1}{2}$}
\matchItem{Б}{$0$}
\matchItem{В}{$\dfrac{1}{2}$}
\matchItem{Г}{$\dfrac{2}{3}$}
\matchItem{Д}{$1$}
}{
\matchingGrid
}
}
% Відповідь: 1--А,\ \ 2--Б,\ \ 3--В
\par\penalty-20
\end{samepage}

\begin{samepage}
% НМТ 2026, сесія 1.06, завдання №16 --- також у темі: Логарифмічні вирази
\zadtask{Узгодьте вираз \mbox{(1--3)} із його значенням \mbox{(А--Д)}, якщо $a = 4$. \nmtyear{2026}

\nmtnobreak
\nopagebreak\nbvspace{0.3cm}
\noindent
\matchingLayout{
\matchHead{Вираз}
\matchItem{1}{$\dfrac{2a^2 - 12a}{a^2 - 36}$}
\matchItem{2}{$0{,}04^{\,\log_5 a}$}
\matchItem{3}{$a^0 - a^{\frac{3}{2}}$}
}{
\matchHead{Значення}
\matchItem{А}{$\dfrac{1}{16}$}
\matchItem{Б}{$-8$}
\matchItem{В}{$-7$}
\matchItem{Г}{$0{,}8$}
\matchItem{Д}{$-\dfrac{1}{16}$}
}{
\matchingGrid
}
}
% Відповідь: 1~--~Г, 2~--~А, 3~--~В
\ifshowsolutions\solution{При $a=4$: \textbf{1)} $\dfrac{2a^2-12a}{a^2-36}=\dfrac{2a(a-6)}{(a-6)(a+6)}=\dfrac{2a}{a+6}=\dfrac{8}{10}=0{,}8$ (Г); \textbf{2)} $0{,}04^{\log_5 4}=(5^{-2})^{\log_5 4}=4^{-2}=\dfrac{1}{16}$ (А); \textbf{3)} $a^0-a^{3/2}=1-8=-7$ (В). Відповідь: \textbf{1\,--\,Г, 2\,--\,А, 3\,--\,В}.}\fi
\par\penalty-20
\end{samepage}

\begin{samepage}
% НМТ 2026, сесія 30.05, завдання №18 --- також у темі: Логарифмічні вирази
\zadtask{Установіть відповідність між виразом \mbox{(1--3)} і його значенням \mbox{(А--Д)}, якщо $a = \dfrac{3}{8}$. \nmtyear{2026}

\nmtnobreak
\nopagebreak\nbvspace{0.2cm}
\noindent
\matchingLayout{
\matchHead{Вираз}
\matchItem{1}{$\sqrt[3]{\dfrac{a}{3}}$}
\matchItem{2}{$\log_2 a + \log_2 \dfrac{16}{3}$}
\matchItem{3}{$\sin(4\pi a)$}
}{
\matchHead{Значення}
\matchItem{А}{$2$}
\matchItem{Б}{$1$}
\matchItem{В}{$\dfrac{1}{2}$}
\matchItem{Г}{$0$}
\matchItem{Д}{$-1$}
}{
\matchingGrid
}
}
% Відповідь: 1~--~В, 2~--~Б, 3~--~Д
\ifshowsolutions\solution{$a=\dfrac38$. \textbf{1}: $\sqrt[3]{\dfrac{a}{3}}=\sqrt[3]{\dfrac18}=\dfrac12$ (В). \textbf{2}: $\log_2\!\left(a\cdot\dfrac{16}{3}\right)=\log_2 2=1$ (Б). \textbf{3}: $\sin\!\left(4\pi\cdot\dfrac38\right)=\sin\dfrac{3\pi}{2}=-1$ (Д). Відповідь: \textbf{1~--~В, 2~--~Б, 3~--~Д}.}\fi
\par\penalty-20
\end{samepage}

\begin{samepage}
% НМТ 2026, сесія 4.06, завдання №16 --- також у темі: Показникова функція і показникові вирази, Логарифмічні вирази
\zadtask{Узгодьте вираз \mbox{(1--3)} із його значенням \mbox{(А--Д)}, якщо $a = \dfrac{2}{3}$. \nmtyear{2026}

\nmtnobreak
\nopagebreak\nbvspace{0.3cm}
\noindent
\matchingLayout{
\matchHead{Вираз}
\matchItem{1}{$6a$}
\matchItem{2}{$\sqrt{\left(2^a\right)^3}$}
\matchItem{3}{$\log_3 (1 - a)$}
}{
\matchHead{Значення виразу}
\matchItem{А}{$4$}
\matchItem{Б}{$-1$}
\matchItem{В}{$1$}
\matchItem{Г}{$\dfrac{20}{3}$}
\matchItem{Д}{$2$}
}{
\matchingGrid
}
}
% Відповідь: 1~--~А, 2~--~Д, 3~--~Б
\ifshowsolutions\solution{При $a=\dfrac{2}{3}$: $6a=4$ (А); $\sqrt{(2^a)^3}=\sqrt{2^{3a}}=\sqrt{2^2}=2$ (Д); $\log_3(1-a)=\log_3\dfrac{1}{3}=-1$ (Б). Відповідь: \textbf{1~--~А, 2~--~Д, 3~--~Б}.}\fi
\par\penalty-20
\end{samepage}

\begin{samepage}
\needspace{0.45\textheight}\typeTitle{Число, задане формулою (золотий перетин $\varphi$, срібний перетин $\delta$)}
% Завдання 14 (на відповідність - золотий перетин)
\zadtask{Число $\varphi = \dfrac{\sqrt{5} + 1}{2}$ називають золотим перетином, що пов'язано з числами Фібоначі. Установіть відповідність між виразом \mbox{(1--3)} та твердженням про його значення \mbox{(А--Д)}, яке є правильним. \nmtyear{2024}}

\nmtnobreak
\nopagebreak\nbvspace{0.3cm}
\noindent
\matchingLayout{
\matchHead{Вираз}
\matchItem{1}{$\varphi \cdot \dfrac{\sqrt{5} - 1}{2}$}

\nmtnobreak
\matchItem{2}{$\log_5(2\varphi - \sqrt{5})$}
\matchItem{3}{$\varphi - 2$}
}{
\matchHead{Твердження про значення виразу}
\matchItem{А}{є натуральним числом}
\matchItem{Б}{є цілим від'ємним числом}
\matchItem{В}{дорівнює 0}
\matchItem{Г}{є раціональним нецілим числом}
\matchItem{Д}{є ірраціональним числом}
}{
\matchingGrid
}
\par\penalty-20\vspace{0.25cm}
\end{samepage}

\begin{samepage}
% Завдання 48 (на відповідність з срібним перетином)
\zadtask{Число $\delta = \sqrt{2} + 1$ називають срібним перетином. Узгодьте вираз \mbox{(1--3)} та твердження про його значення \mbox{(А--Д)}, яке є правильним для цього виразу. \nmtyear{2025}}

\nmtnobreak
\nopagebreak\nbvspace{0.3cm}
\matchingLayout{
\matchHead{Вираз}
\matchItem{1}{$\dfrac{1}{\delta}$}
\matchItem{2}{$\log_2(\delta - 1)$}
\matchItem{3}{$\sqrt{2} - \delta$}
}{
\matchHead{Твердження про значення виразу}
\matchItem{А}{є натуральним числом}
\matchItem{Б}{є цілим від'ємним числом}
\matchItem{В}{є раціональним нецілим числом}
\matchItem{Г}{є ірраціональним додатним числом}
\matchItem{Д}{є ірраціональним від'ємним числом}
}{
\matchingGrid
}
\par\penalty-20\vspace{0.25cm}
\end{samepage}

\begin{samepage}
% Завдання 32 (на відповідність - золотий перетин з проміжками)
\zadtask{Число $\varphi = \dfrac{\sqrt{5} + 1}{2}$ називають золотим перетином, яке тісно пов'язане з послідовністю чисел Фібоначчі. Узгодьте вираз \mbox{(1--3)} та проміжок \mbox{(А--Д)}, якому належить значення цього виразу. \nmtyear{2025}}

\nmtnobreak
\nopagebreak\nbvspace{0.3cm}
\noindent
\matchingLayout{
\matchHead{Вираз}
\matchItem{1}{$\varphi^0$}
\matchItem{2}{$(\sqrt{5} - 1)\varphi$}
\matchItem{3}{$\log_{\frac{1}{5}}(2\varphi - 1)$}
}{
\matchHead{Проміжок}
\matchItem{А}{$(-\infty; -1]$}
\matchItem{Б}{$(-1; 0]$}
\matchItem{В}{$(0; 1]$}
\matchItem{Г}{$(1; 2]$}
\matchItem{Д}{$(2; +\infty)$}
}{
\matchingGrid
}
\par\penalty-20\vspace{0.25cm}
\end{samepage}

\begin{samepage}
% Завдання 25 (на відповідність з числовою прямою - золотий перетин)
\zadtask{Число $\varphi = \dfrac{\sqrt{5} + 1}{2}$ називають золотим перетином, яке тісно пов'язане з послідовністю чисел Фібоначчі. Узгодьте вираз \mbox{(1--3)} й точку \mbox{(А--Д)} на координатній прямій (див. рисунок), координатою якої є значення виразу. \nmtyear{2025}}

\nmtnobreak
\nopagebreak\nbvspace{0.3cm}
\begin{center}
\begin{nmtfit}\begin{tikzpicture}[scale=1.3]
    \draw[->] (-3,0) -- (3,0);
    \foreach \x/\name in {-2/K, -1/L, 0/M, 1/N, 2/P} {
        \draw (\x,0.1) -- (\x,-0.1);
        \node[above] at (\x,0.15) {$\name$};
        \node[below] at (\x,-0.15) {$\x$};
    }
\end{tikzpicture}\end{nmtfit}
\end{center}

\nmtnobreak
\nopagebreak\nbvspace{0.3cm}
\noindent
\matchingLayout{
\matchHead{Вираз}
\matchItem{1}{$\dfrac{\sqrt{5} - 1}{2} - \varphi$}

\nmtnobreak
\matchItem{2}{$(1 - \sqrt{5}) \cdot \varphi$}
\matchItem{3}{$\log_5(2\varphi - \sqrt{5})$}
}{
\matchHead{Точка}
\matchItem{А}{$K$}
\matchItem{Б}{$L$}
\matchItem{В}{$M$}
\matchItem{Г}{$N$}
\matchItem{Д}{$P$}
}{
\matchingGrid
}
\par\penalty-20\vspace{0.25cm}
\end{samepage}

\begin{samepage}
\needspace{0.45\textheight}\typeTitle{Проміжок, якому належить значення числового виразу}
% Завдання 403 (НМТ 2022, демоваріант) --- з тематичного збірника «Числа і вирази»
\zadtask{Установіть відповідність між виразом \mbox{(1--3)} та проміжком \mbox{(А--Д)}, якому належить його значення. \nmtyear{2022}}

\nmtnobreak
\nopagebreak\nbvspace{0.3cm}
\matchingLayout{
\matchHead{Вираз}
\matchItem{1}{$|-1{,}6| + 2$}
\matchItem{2}{$\dfrac{\sqrt{24}}{\sqrt{3}}$}

\nmtnobreak
\matchItem{3}{$2\cos\dfrac{\pi}{3}$}
}{
\matchHead{Проміжок}
\matchItem{А}{$(-\infty; 0)$}
\matchItem{Б}{$[0; 1)$}
\matchItem{В}{$[1; 2)$}
\matchItem{Г}{$[2; 3)$}
\matchItem{Д}{$[3; +\infty)$}
}{
\matchingGrid
}
\par\penalty-20\vspace{0.25cm}
\end{samepage}

\begin{samepage}
% Завдання 18 (на відповідність)
\zadtask{Установіть відповідність між виразом \mbox{(1--3)} і проміжком \mbox{(А--Д)}, якому належить значення цього виразу. \nmtyear{2023}}

\nmtnobreak
\nopagebreak\nbvspace{0.3cm}
\matchingLayout{
\matchHead{Вираз}
\matchItem{1}{$\sqrt{(-3)^2}$}
\matchItem{2}{$\sqrt{3} - 1$}
\matchItem{3}{$\log_{\frac{1}{3}} \sqrt{3}$}
}{
\matchHead{Проміжок}
\matchItem{А}{$[-3; -1)$}
\matchItem{Б}{$[-1; 0)$}
\matchItem{В}{$[0; 1)$}
\matchItem{Г}{$[1; 2)$}
\matchItem{Д}{$[2; 3]$}
}{
\matchingGrid
}
\par\penalty-20\vspace{0.25cm}
\end{samepage}

\begin{samepage}
% Завдання 7 (на відповідність)
\zadtask{Установіть відповідність між виразом \mbox{(1--3)} та проміжком \mbox{(А--Д)}, якому належить значення цього виразу. \nmtyear{2024}}

\nmtnobreak
\nopagebreak\nbvspace{0.3cm}
\noindent
\matchingLayout{
\matchHead{Вираз}
\matchItem{1}{$\sqrt{2} \cdot \sqrt{18}$}
\matchItem{2}{$|\sqrt{2} - 2|$}
\matchItem{3}{$\log_{\sqrt{2}} 0{,}5$}
}{
\matchHead{Проміжок}
\matchItem{А}{$(-\infty; -2)$}
\matchItem{Б}{$[-2; 0)$}
\matchItem{В}{$[0; 1)$}
\matchItem{Г}{$[1; 2)$}
\matchItem{Д}{$[2; +\infty)$}
}{
\matchingGrid
}
\par\penalty-20\vspace{0.25cm}
\end{samepage}

\begin{samepage}
% Завдання 26 (на відповідність)
\zadtask{Установіть відповідність між виразом \mbox{(1--3)} та проміжком \mbox{(А--Д)}, якому належить значення цього виразу. \nmtyear{2024}}

\nmtnobreak
\nopagebreak\nbvspace{0.3cm}
\matchingLayout{
\matchHead{Вираз}
\matchItem{1}{$(-\sqrt{2})^2$}
\matchItem{2}{$1 - \sqrt{2}$}
\matchItem{3}{$\left(\dfrac{1}{3}\right)^{\log_3 \sqrt{2}}$}
}{
\matchHead{Проміжок}
\matchItem{А}{$[-4; -1)$}
\matchItem{Б}{$[-1; 0)$}
\matchItem{В}{$[0; 1)$}
\matchItem{Г}{$[1; 2)$}
\matchItem{Д}{$[2; 5)$}
}{
\matchingGrid
}
\par\penalty-20\vspace{0.25cm}
\end{samepage}

\begin{samepage}
% НМТ 2026, сесія 8.06, завдання №16 --- основна тема: Логарифмічні вирази
\zadtask{Узгодьте вираз \mbox{(1--3)} із проміжком \mbox{(А--Д)}, якому належить значення цього виразу. \nmtyear{2026}\par
\nopagebreak\nbvspace{0.3cm}
\noindent
\matchingLayout{
\matchHead{Вираз}
\matchItem{1}{$\log_{0{,}5} 2$}
\matchItem{2}{$\sqrt{3^2 + (-2)^2}$}
\matchItem{3}{$\left| 2 \cdot \dfrac{2}{7} - 2\dfrac{2}{7} \right|$}
}{
\matchHead{Проміжок}
\matchItem{А}{$[-2; -1)$}
\matchItem{Б}{$[-1; 1)$}
\matchItem{В}{$[1; 2)$}
\matchItem{Г}{$[2; 3)$}
\matchItem{Д}{$[3; 4)$}
}{
\matchingGrid
}
}
\par\penalty-20
\end{samepage}

\begin{samepage}
% НМТ 2026, сесія 9.06, завдання №16 --- також у темі: Логарифмічні вирази
\zadtask{Узгодьте вираз \mbox{(1--3)} із проміжком \mbox{(А--Д)}, якому належить значення цього виразу. \nmtyear{2026}\par
\nopagebreak\nbvspace{0.3cm}
\noindent
\matchingLayout{
\matchHead{Вираз}
\matchItem{1}{$\dfrac{\sqrt{200}}{\sqrt{8}}$}
\matchItem{2}{$\log_8 200$}
\matchItem{3}{$8 \cos 120^\circ$}
}{
\matchHead{Проміжок}
\matchItem{А}{$[-5; 0)$}
\matchItem{Б}{$[0; 2)$}
\matchItem{В}{$[2; 4)$}
\matchItem{Г}{$[4; 6)$}
\matchItem{Д}{$[6; +\infty)$}
}{
\matchingGrid
}
}
\par\penalty-20
\end{samepage}

\begin{samepage}
% НМТ 2026, сесія 2.06, завдання №16 --- також у темі: Логарифмічні вирази
\zadtask{Узгодьте вираз \mbox{(1--3)} із проміжком \mbox{(А--Д)}, якому належить значення цього виразу. \nmtyear{2026}

\nmtnobreak
\nopagebreak\nbvspace{0.3cm}
\noindent
\matchingLayout{
\matchHead{Вираз}
\matchItem{1}{$\dfrac{7}{5} - \dfrac{5}{3}$}
\matchItem{2}{$\sqrt{3} \cdot \sqrt{5}$}
\matchItem{3}{$\log_4 8$}
}{
\matchHead{Проміжок}
\matchItem{А}{$[-1;\, 0)$}
\matchItem{Б}{$[0;\, 1)$}
\matchItem{В}{$[1;\, 2)$}
\matchItem{Г}{$[2;\, 4)$}
\matchItem{Д}{$[4;\, +\infty)$}
}{
\matchingGrid
}
}
% Відповідь: 1~--~А, 2~--~Г, 3~--~В
\ifshowsolutions\solution{$\dfrac{7}{5}-\dfrac{5}{3}=\dfrac{21-25}{15}=-\dfrac{4}{15}\in[-1;0)$ -- \textbf{А}; $\sqrt{3}\cdot\sqrt{5}=\sqrt{15}\approx3{,}87\in[2;4)$ -- \textbf{Г}; $\log_4 8=\dfrac{3}{2}=1{,}5\in[1;2)$ -- \textbf{В}. Відповідь: \textbf{1--А, 2--Г, 3--В}.}\fi
\par\penalty-20
\end{samepage}

\begin{samepage}
\needspace{0.45\textheight}\typeTitle{Проміжок: вирази з числом $\pi$ або $e$}
\noindent\zadnum Установіть відповідність між виразом \mbox{(1--3)} і проміжком \mbox{(А--Д)}, якому належить значення цього виразу, якщо $e \approx 2{,}7$. \nmtyear{2023}\par
\nopagebreak\nbvspace{0.3cm}

\nmtnobreak
\matchingLayout{
\matchHead{Вираз}

\nmtnobreak
\matchItem{1}{$\ln e$}
\matchItem{2}{$e - 4$}
\matchItem{3}{$e^2$}
}{
\matchHead{Проміжок}

\nmtnobreak
\matchItem{А}{$(-\infty; -2]$}
\matchItem{Б}{$(-2; 0]$}
\matchItem{В}{$(0; 2]$}
\matchItem{Г}{$(2; 6]$}
\matchItem{Д}{$(6; +\infty)$}
}{
\matchingGrid
}
\par\penalty-20\vspace{0.25cm}
\end{samepage}

\begin{samepage}
% === ЗАВДАННЯ 10 (Відповідність) ===
\noindent\zadnum Установіть відповідність між виразом \mbox{(1--3)} і проміжком \mbox{(А--Д)}, якому належить значення цього виразу. \nmtyear{2023}

\nmtnobreak
\nopagebreak\nbvspace{0.3cm}

\nmtnobreak
\matchingLayout{
\matchHead{Вираз}
    \matchItem{1}{$\dfrac{\pi}{2}$}
\matchItem{2}{\trigStyle{$\sin \pi$}}
\matchItem{3}{$\log_{\pi} \dfrac{1}{\pi}$}
}{
\matchHead{Проміжок}
    \matchItem{А}{$[-2; -1)$}
\matchItem{Б}{$[-1; 0)$}
\matchItem{В}{$[0; 1)$}
\matchItem{Г}{$[1; 2)$}
\matchItem{Д}{$[2; 3)$}
}{
    \matchingGrid
}
\par\penalty-20\vspace{0.25cm}
\end{samepage}

\begin{samepage}
% Завдання 3 (на відповідність)
\zadtask{Установіть відповідність між виразом \mbox{(1--3)} і проміжком \mbox{(А--Д)}, якому належить значення цього виразу. \nmtyear{2023}}

\nmtnobreak
\nopagebreak\nbvspace{0.3cm}
\noindent
\matchingLayout{
\matchHead{Вираз}
\matchItem{1}{$\cos \dfrac{\pi}{2}$}
\matchItem{2}{$|\pi - 5|$}
\matchItem{3}{$2^\pi$}
}{
\matchHead{Проміжок}
\matchItem{А}{$(-\infty; 0]$}
\matchItem{Б}{$(0; 2)$}
\matchItem{В}{$[2; 4)$}
\matchItem{Г}{$[4; 8)$}
\matchItem{Д}{$[8; +\infty)$}
}{
\matchingGrid
}
\par\penalty-20\vspace{0.25cm}
\end{samepage}

\begin{samepage}
% Завдання 9 (на відповідність)
\zadtask{Установіть відповідність між виразом \mbox{(1--3)} і проміжком \mbox{(А--Д)}, якому належить значення цього виразу. \nmtyear{2023}}

\nmtnobreak
\nopagebreak\nbvspace{0.3cm}
\matchingLayout{
\matchHead{Вираз}
\matchItem{1}{$\ln \dfrac{1}{e}$}
\matchItem{2}{$|e - 5|$}
\matchItem{3}{$e^0$}
}{
\matchHead{Проміжок}
\matchItem{А}{$(-\infty; -1)$}
\matchItem{Б}{$[-1; 0)$}
\matchItem{В}{$[0; 1)$}
\matchItem{Г}{$[1; 2)$}
\matchItem{Д}{$(2; +\infty)$}
}{
\matchingGrid
}
\par\penalty-20\vspace{0.25cm}
\end{samepage}

\begin{samepage}
% Завдання 9 (на відповідність)
\zadtask{Установіть відповідність між виразом \mbox{(1--3)} та проміжком \mbox{(А--Д)}, якому належить значення цього виразу. \nmtyear{2024}}

\nmtnobreak
\nopagebreak\nbvspace{0.3cm}
\noindent
\matchingLayout{
\matchHead{Вираз}
\matchItem{1}{$\mathrm{tg}\,\dfrac{\pi}{3}$}
\matchItem{2}{$1 - \pi$}
\matchItem{3}{$\left(\dfrac{1}{2}\right)^{\log_2 \pi}$}
}{
\matchHead{Проміжок}
\matchItem{А}{$[-5; -2)$}
\matchItem{Б}{$[-2; 0)$}
\matchItem{В}{$[0; 1)$}
\matchItem{Г}{$[1; 2)$}
\matchItem{Д}{$[2; 5)$}
}{
\matchingGrid
}
\par\penalty-20\vspace{0.25cm}
\end{samepage}

\begin{samepage}
% Завдання 3 (на відповідність)
\zadtask{Установіть відповідність між виразом \mbox{(1--3)} та проміжком \mbox{(А--Д)}, якому належить значення цього виразу. \nmtyear{2024}}

\nmtnobreak
\nopagebreak\nbvspace{0.3cm}
\noindent
\matchingLayout{
\matchHead{Вираз}
\matchItem{1}{$\cos \dfrac{\pi}{3}$}
\matchItem{2}{$2\pi - 5$}
\matchItem{3}{$\log_3 \pi - \log_3 (3\pi)$}
}{
\matchHead{Проміжок}
\matchItem{А}{$[-4; -1)$}
\matchItem{Б}{$[-1; 0)$}
\matchItem{В}{$[0; 1)$}
\matchItem{Г}{$[1; 2)$}
\matchItem{Д}{$[2; 5)$}
}{
\matchingGrid
}
\par\penalty-20\vspace{0.25cm}
\end{samepage}

\begin{samepage}
% Завдання 9 (на відповідність)
\zadtask{Установіть відповідність між виразом \mbox{(1--3)}, де $\pi$ --- відома математична константа, та проміжком \mbox{(А--Д)}, якому належить його значення. \nmtyear{2025}}

\nmtnobreak
\nopagebreak\nbvspace{0.3cm}
\noindent
\matchingLayout{
\matchHead{Вираз}
\matchItem{1}{$\cos \dfrac{\pi}{3}$}
\matchItem{2}{$\pi - 4$}
\matchItem{3}{$4^{\log_4 \pi}$}
}{
\matchHead{Проміжок}
\matchItem{А}{$[-2; -1)$}
\matchItem{Б}{$[-1; 0)$}
\matchItem{В}{$[0; 1)$}
\matchItem{Г}{$[1; 2)$}
\matchItem{Д}{$[2; 4)$}
}{
\matchingGrid
}
\par\penalty-20\vspace{0.25cm}
\end{samepage}

\begin{samepage}
\needspace{0.45\textheight}\typeTitle{Точка на координатній прямій, координата якої --- значення виразу}
% Завдання 5 (на відповідність з числовою прямою)
\zadtask{Установіть відповідність між виразом \mbox{(1--3)} та точкою \mbox{(А--Д)} на координатній прямій (див. рисунок), координатою якої є значення цього виразу за $a = 0{,}5$. \nmtyear{2023}}

\nmtnobreak
\nopagebreak\nbvspace{0.3cm}
\begin{center}
\begin{nmtfit}\begin{tikzpicture}[scale=1.3]
    \draw[->] (-3,0) -- (3,0);
    \foreach \x/\name in {-2/K, -1/L, 0/M, 1/N, 2/P} {
        \draw (\x,0.1) -- (\x,-0.1);
        \node[above] at (\x,0.15) {$\name$};
        \node[below] at (\x,-0.15) {$\x$};
    }
\end{tikzpicture}\end{nmtfit}
\end{center}

\nmtnobreak
\nopagebreak\nbvspace{0.3cm}
\matchingLayout{
\matchHead{Вираз}
\matchItem{1}{$|a - 2{,}5|$}
\matchItem{2}{$a^0$}
\matchItem{3}{$\log_2 a$}
}{
\matchHead{Точка}
\matchItem{А}{$K$}
\matchItem{Б}{$L$}
\matchItem{В}{$M$}
\matchItem{Г}{$N$}
\matchItem{Д}{$P$}
}{
\matchingGrid
}
\par\penalty-20\vspace{0.25cm}
\end{samepage}

\begin{samepage}
% Завдання 19 (на відповідність з числовою прямою)
\zadtask{Узгодьте вираз \mbox{(1--3)} з точкою \mbox{(А--Д)} на координатній прямій, координатою якої є значення виразу. \nmtyear{2024}}

\nmtnobreak
\nopagebreak\nbvspace{0.3cm}
\begin{center}
\begin{nmtfit}\begin{tikzpicture}[scale=1.3]
    \draw[->] (-3,0) -- (3,0);
    \foreach \x/\name in {-2/K, -1/L, 0/M, 1/N, 2/P} {
        \draw (\x,0.1) -- (\x,-0.1);
        \node[above] at (\x,0.15) {$\name$};
        \node[below] at (\x,-0.15) {$\x$};
    }
\end{tikzpicture}\end{nmtfit}
\end{center}

\nmtnobreak
\nopagebreak\nbvspace{0.3cm}
\matchingLayout{
\matchHead{Вираз}
\matchItem{1}{$2\pi \cdot \pi^{-1}$}
\matchItem{2}{$\mathrm{tg}\,\dfrac{5\pi}{4}$}
\matchItem{3}{$\log_\pi \dfrac{1}{\pi^2}$}
}{
\matchHead{Точка}
\matchItem{А}{$K$}
\matchItem{Б}{$L$}
\matchItem{В}{$M$}
\matchItem{Г}{$N$}
\matchItem{Д}{$P$}
}{
\matchingGrid
}
\par\penalty-20\vspace{0.25cm}
\end{samepage}

\begin{samepage}
\noindent\zadnum Узгодьте вираз \mbox{(1--3)} з точкою \mbox{(А--Д)} на координатній прямій (див. рисунок), координатою якої є значення цього виразу. \nmtyear{2024}

\nmtnobreak
\begin{center}
    \begin{nmtfit}\begin{tikzpicture}[scale=1.2]
        \draw[->, thick] (-2.5,0) -- (2.5,0);
        \foreach \x/\l in {-2/K, -1/L, 0/M, 1/N, 2/P} {
            \path (\x,0) node[nmt dot] {} node[above] {$\l$};
            \node[below] at (\x,0) {$\x$};
        }
    \end{tikzpicture}\end{nmtfit}
\end{center}

\nmtnobreak
\nopagebreak\nbvspace{0.3cm}

\nmtnobreak
\matchingLayout{
\matchHead{Вираз}
    \matchItem{1}{\trigStyle{$\sin^2 \dfrac{\pi}{6} + \cos^2 \dfrac{\pi}{6}$}}
\matchItem{2}{$\dfrac{\pi^2-4}{\pi-2} - \pi$}
\matchItem{3}{$\log_3 \pi^0$}
}{
\matchHead{Точка}
    \matchItem{А}{$K$}
\matchItem{Б}{$L$}
\matchItem{В}{$M$}
\matchItem{Г}{$N$}
\matchItem{Д}{$P$}
}{
    \matchingGrid
}
\par\penalty-20\vspace{0.25cm}
\end{samepage}

\begin{samepage}
% Завдання 14 (на відповідність з числовою прямою)
\zadtask{Узгодьте вираз \mbox{(1--3)} з точкою \mbox{(А--Д)} на координатній прямій, координатою якої є значення виразу, якщо $a = -2$. \nmtyear{2024}}

\nmtnobreak
\nopagebreak\nbvspace{0.3cm}
\begin{center}
\begin{nmtfit}\begin{tikzpicture}[scale=1.3]
    \draw[->] (-3,0) -- (3,0);
    \foreach \x/\name in {-2/K, -1/L, 0/M, 1/N, 2/P} {
        \draw (\x,0.1) -- (\x,-0.1);
        \node[above] at (\x,0.15) {$\name$};
        \node[below] at (\x,-0.15) {$\x$};
    }
\end{tikzpicture}\end{nmtfit}
\end{center}

\nmtnobreak
\nopagebreak\nbvspace{0.3cm}
\noindent
\matchingLayout{
\matchHead{Вираз}
\matchItem{1}{$|a|$}
\matchItem{2}{$a^0$}
\matchItem{3}{$\mathrm{tg}\,(\pi a)$}
}{
\matchHead{Точка}
\matchItem{А}{$K$}
\matchItem{Б}{$L$}
\matchItem{В}{$M$}
\matchItem{Г}{$N$}
\matchItem{Д}{$P$}
}{
\matchingGrid
}
\par\penalty-20\vspace{0.25cm}
\end{samepage}

\begin{samepage}
% Завдання 22 (на відповідність з числовою прямою)
\zadtask{Узгодьте вираз \mbox{(1--3)} з точкою \mbox{(А--Д)} на координатній прямій, координатою якої є значення виразу. \nmtyear{2024}}

\nmtnobreak
\nopagebreak\nbvspace{0.3cm}
\begin{center}
\begin{nmtfit}\begin{tikzpicture}[scale=1.3]
    \draw[->] (-3,0) -- (3,0);
    \foreach \x/\name in {-2.414/K, -1/L, 0/M, 1/N, 2.414/P} {
        \draw (\x,0.1) -- (\x,-0.1);
        \node[above] at (\x,0.15) {$\name$};
    }
    \foreach \x in {-1, 0, 1} \node[below] at (\x,-0.15) {$\x$};
\end{tikzpicture}\end{nmtfit}
\end{center}

\nmtnobreak
\nopagebreak\nbvspace{0.3cm}
\matchingLayout{
\matchHead{Вираз}
\matchItem{1}{$\log_{\sqrt{2}} \cos 360°$}

\nmtnobreak
\matchItem{2}{$\dfrac{1}{\sqrt{2} - 1}$}

\nmtnobreak
\matchItem{3}{$1 - (\sqrt{2})^2$}
}{
\matchHead{Точка}
\matchItem{А}{$K$}
\matchItem{Б}{$L$}
\matchItem{В}{$M$}
\matchItem{Г}{$N$}
\matchItem{Д}{$P$}
}{
\matchingGrid
}
\par\penalty-20\vspace{0.25cm}
\end{samepage}

\begin{samepage}
% Завдання 2 (на відповідність з числовою прямою)
\zadtask{Узгодьте вираз \mbox{(1--3)} й точку \mbox{(А--Д)} на координатній прямій (див. рисунок), координатою якої є значення виразу, де $e \approx 2{,}7$ --- основа натурального логарифма (число Ейлера). \nmtyear{2025}}

\nmtnobreak
\nopagebreak\nbvspace{0.3cm}
\begin{center}
\begin{nmtfit}\begin{tikzpicture}[scale=1.3]
    \draw[->] (-3,0) -- (3,0);
    \foreach \x/\name in {-2/K, -1/L, 0/M, 1/N, 2/P} {
        \draw (\x,0.1) -- (\x,-0.1);
        \node[above] at (\x,0.15) {$\name$};
        \node[below] at (\x,-0.15) {$\x$};
    }
\end{tikzpicture}\end{nmtfit}
\end{center}

\nmtnobreak
\nopagebreak\nbvspace{0.3cm}
\noindent
\matchingLayout{
\matchHead{Вираз}
\matchItem{1}{$2e \cdot \dfrac{1}{e}$}
\matchItem{2}{$\ln 1$}
\matchItem{3}{$(e - 1)(e + 1) - e^2$}
}{
\matchHead{Точка}
\matchItem{А}{$K$}
\matchItem{Б}{$L$}
\matchItem{В}{$M$}
\matchItem{Г}{$N$}
\matchItem{Д}{$P$}
}{
\matchingGrid
}
\par\penalty-20\vspace{0.25cm}
\end{samepage}

\begin{samepage}
\noindent\zadnum Узгодьте вираз \mbox{(1--3)} й точку \mbox{(А--Д)} на координатній прямій (див. рисунок), координатою якої є значення виразу, де $e \approx 2{,}7$ -- основа натурального логарифма (число Ейлера). \nmtyear{2025}\par
\nopagebreak\nbvspace{0.3cm}

\nmtnobreak
\begin{center}
\begin{nmtfit}\begin{tikzpicture}
\draw[thick, ->] (-2.8,0) -- (2.8,0);
\foreach \x/\label in {-2/K, -1/L, 0/M, 1/N, 2/P} {
    \draw[thick] (\x,0.1) -- (\x,-0.1) node[below] {\normalsize $\label$};
}
\node[above] at (-2,0.1) {\normalsize $-1$};
\node[above] at (0,0.1) {\normalsize $0$};
\node[above] at (2,0.1) {\normalsize $1$};
\end{tikzpicture}\end{nmtfit}
\end{center}

\nmtnobreak
\nopagebreak\nbvspace{0.3cm}

\nmtnobreak
\matchingLayout{\textit{Вираз}

\nmtnobreak
\matchItem{1}{$e \cdot \displaystyle\frac{1}{2e}$}
\matchItem{2}{$\ln 2e - \ln 2$}
\matchItem{3}{$|1 - e| - e$}
}{
\matchHead{Точка}

\nmtnobreak
\matchItem{А}{$K$}
\matchItem{Б}{$L$}
\matchItem{В}{$M$}
\matchItem{Г}{$N$}
\matchItem{Д}{$P$}
}{
\matchingGrid
}
\par\penalty-20\vspace{0.25cm}
\end{samepage}

\begin{samepage}
\needspace{0.45\textheight}\typeTitle{Твердження про значення виразу:\\ натуральне, ціле, раціональне, ірраціональне}
\noindent\zadnum Установіть відповідність між виразом (1–3) і множиною (А–Д), якій належить значення цього виразу, якщо $a = 5$. \nmtyear{2023}\par
\nopagebreak\nbvspace{0.3cm}

\nmtnobreak
\noindent
\matchingLayout{
\matchHead{Вираз}
    \matchItem{1}{$\log_{0{,}2} a$}
\matchItem{2}{$2^{-a}$}
\matchItem{3}{$\displaystyle\frac{10}{a}$}
}{
\matchHead{Множина}
    \matchItem{А}{ірраціональних додатних чисел}
\matchItem{Б}{натуральних чисел}
\matchItem{В}{цілих від'ємних чисел}
\matchItem{Г}{раціональних нецілих чисел}
\matchItem{Д}{ірраціональних від'ємних чисел}
}{
\matchingGrid
}
\par\penalty-20\vspace{0.25cm}
\end{samepage}

\begin{samepage}
% Завдання 23 (на відповідність)
\zadtask{Установіть відповідність між виразом \mbox{(1--3)} та твердженням про його значення \mbox{(А--Д)}, яке є правильним. \nmtyear{2024}}

\nmtnobreak
\nopagebreak\nbvspace{0.3cm}
\matchingLayout{
\matchHead{Вираз}
\matchItem{1}{$\log_\pi 1$}
\matchItem{2}{$\sin\left(-\dfrac{\pi}{6}\right)$}
\matchItem{3}{$\pi^3 \cdot \pi^{-4}$}
}{
\matchHead{Твердження про значення виразу}
\matchItem{А}{є нецілим додатним числом}
\matchItem{Б}{є нецілим від'ємним числом}
\matchItem{В}{дорівнює 0}
\matchItem{Г}{є цілим додатним числом}
\matchItem{Д}{є цілим від'ємним числом}
}{
\matchingGrid
}
\par\penalty-20\vspace{0.25cm}
\end{samepage}

\begin{samepage}
% Завдання 22 (на відповідність)
\zadtask{До кожного початку речення \mbox{(1--3)} доберіть його закінчення \mbox{(А--Д)} так, щоб утворилося правильне твердження, якщо $a = 3$. \nmtyear{2024}}

\nmtnobreak
\nopagebreak\nbvspace{0.3cm}
\matchingLayout{
\matchHead{Початок речення}
\matchItem{1}{Значення виразу $a^{-1}$}
\matchItem{2}{Значення виразу $a^0$}
\matchItem{3}{Значення виразу $\sin(\pi a)$}
}{
\matchHead{Закінчення речення}
\matchItem{А}{є раціональним нецілим числом.}
\matchItem{Б}{є ірраціональним числом.}
\matchItem{В}{є натуральним числом.}
\matchItem{Г}{дорівнює нулю.}
\matchItem{Д}{є цілим від'ємним числом.}
}{
\matchingGrid
}
\par\penalty-20\vspace{0.25cm}
\end{samepage}

\begin{samepage}
% Завдання 8 (на відповідність)
\zadtask{Установіть відповідність між виразом \mbox{(1--3)} та твердженням про його значення \mbox{(А--Д)}, яке є правильним. \nmtyear{2024}}

\nmtnobreak
\nopagebreak\nbvspace{0.3cm}
\noindent
\matchingLayout{
\matchHead{Вираз}
\matchItem{1}{$\sin \dfrac{7\pi}{2}$}
\matchItem{2}{$\pi^{\cos 90°}$}
\matchItem{3}{$\dfrac{\pi}{3}$}
}{
\matchHead{Твердження про значення виразу}
\matchItem{А}{є раціональним нецілим числом}
\matchItem{Б}{є ірраціональним числом}
\matchItem{В}{дорівнює $0$}
\matchItem{Г}{є натуральним числом}
\matchItem{Д}{є цілим від'ємним числом}
}{
\matchingGrid
}
\par\penalty-20\vspace{0.25cm}
\end{samepage}

\begin{samepage}
% Завдання 11 (на відповідність)
\zadtask{Установіть відповідність між виразом \mbox{(1--3)} та твердженням про його значення \mbox{(А--Д)}, яке є правильним, якщо $e \approx 2{,}7$. \nmtyear{2024}}

\nmtnobreak
\nopagebreak\nbvspace{0.3cm}
\noindent
\matchingLayout{
\matchHead{Вираз}
\matchItem{1}{$2e \cdot \dfrac{1}{e}$}
\matchItem{2}{$(e - 1)(e + 1)$}
\matchItem{3}{$\ln\left(\sqrt{e} \cdot e^{-\frac{1}{2}}\right)$}
}{
\matchHead{Твердження про значення виразу}
\matchItem{А}{є простим числом}
\matchItem{Б}{є цілим від'ємним числом}
\matchItem{В}{дорівнює $0$}
\matchItem{Г}{є нецілим додатним числом}
\matchItem{Д}{є нецілим від'ємним числом}
}{
\matchingGrid
}
\par\penalty-20\vspace{0.25cm}
\end{samepage}

\begin{samepage}
% Завдання 13 (на відповідність)
\zadtask{Установіть відповідність між виразом \mbox{(1--3)} та твердженням про його значення \mbox{(А--Д)}, яке є правильним. \nmtyear{2024}}

\nmtnobreak
\nopagebreak\nbvspace{0.3cm}
\noindent
\matchingLayout{
\matchHead{Вираз}
\matchItem{1}{$(\sqrt{2} + 5)(\sqrt{2} - 5)$}
\matchItem{2}{$2\log_2 \sqrt{8}$}
\matchItem{3}{$|1 - \sqrt{2}|$}
}{
\matchHead{Твердження про значення виразу}
\matchItem{А}{є цілим додатним числом}
\matchItem{Б}{є цілим від'ємним числом}
\matchItem{В}{дорівнює $0$}
\matchItem{Г}{є нецілим додатним числом}
\matchItem{Д}{є нецілим від'ємним числом}
}{
\matchingGrid
}
\par\penalty-20\vspace{0.25cm}
\end{samepage}

\begin{samepage}
\noindent\zadnum Установіть відповідність між виразом \mbox{(1--3)} та твердженням про його значення \mbox{(А--Д)}, яке є правильним. \nmtyear{2024}

\nmtnobreak
\nopagebreak\nbvspace{0.3cm}

\nmtnobreak
\matchingLayout{
\matchHead{Вираз}
    \matchItem{1}{\trigStyle{$\cos 2\pi$}}
\matchItem{2}{$\log_{\pi} \dfrac{1}{\pi}$}
\matchItem{3}{$\pi^2 - 9$}
}{
\matchHead{Твердження про значення виразу}
    \matchItem{А}{є цілим додатним числом}
\matchItem{Б}{є цілим від'ємним числом}
\matchItem{В}{дорівнює 0}
\matchItem{Г}{є нецілим додатним числом}
\matchItem{Д}{є нецілим від'ємним числом}
}{
    \matchingGrid
}
\par\penalty-20\vspace{0.25cm}
\end{samepage}

\begin{samepage}
% Завдання 27 (на відповідність)
\zadtask{Узгодьте вираз \mbox{(1--3)} і твердження про його значення \mbox{(А--Д)}, яке є правильним для цього виразу, де $e \approx 2{,}7$ --- основа натурального логарифма (число Ейлера). \nmtyear{2025}}

\nmtnobreak
\nopagebreak\nbvspace{0.3cm}
\noindent
\matchingLayout{
\matchHead{Вираз}
\matchItem{1}{$\ln\dfrac{1}{e}$}
\matchItem{2}{$e + 1$}
\matchItem{3}{$(e - 1)^2 - e^2 + 2e$}
}{
\matchHead{Твердження про значення виразу}
\matchItem{А}{є додатним цілим числом}
\matchItem{Б}{є від'ємним цілим числом}
\matchItem{В}{є додатним нецілим числом}
\matchItem{Г}{є від'ємним нецілим числом}
\matchItem{Д}{дорівнює 0}
}{
\matchingGrid
}
\par\penalty-20\vspace{0.25cm}
\end{samepage}

\begin{samepage}
\noindent\zadnum Узгодьте вираз \mbox{(1--3)} і твердження про його значення \mbox{(А--Д)}, яке є правильним для цього виразу. \nmtyear{2025}

\nmtnobreak
\nopagebreak\nbvspace{0.3cm}

\nmtnobreak
\noindent
\matchingLayout{
\matchHead{Вираз}
    \matchItem{1}{$(3\pi - 1)^0$}
\matchItem{2}{$\log_{\pi} \dfrac{1}{\pi^3}$}
\matchItem{3}{\trigStyle{$\operatorname{tg} \dfrac{\pi}{3}$}}
}{
\matchHead{Твердження про значення виразу}
    \matchItem{А}{є натуральним числом}
\matchItem{Б}{є цілим недодатним числом}
\matchItem{В}{є раціональним нецілим числом}
\matchItem{Г}{є ірраціональним додатним числом}
\matchItem{Д}{є ірраціональним від'ємним числом}
}{
\matchingGrid
}
\par\penalty-20\vspace{0.25cm}
\end{samepage}

\begin{samepage}
\needspace{0.45\textheight}\typeTitle{Тотожно рівний вираз зі змінною}
% Завдання 17 (на відповідність)
\zadtask{До кожного виразу \mbox{(1--3)} доберіть тотожно рівний йому вираз \mbox{(А--Д)}, якщо $a \neq \sqrt{2}$. \nmtyear{2023}}

\nmtnobreak
\nopagebreak\nbvspace{0.3cm}
\matchingLayout{
\matchHead{Вираз}
\matchItem{1}{$\dfrac{a^2 - 2a\sqrt{2} + (\sqrt{2})^2}{a - \sqrt{2}}$}

\nmtnobreak
\matchItem{2}{$\dfrac{a^2 - 2}{a - \sqrt{2}}$}

\nmtnobreak
\matchItem{3}{$\log_b b^a + 3$}
}{
\matchHead{Тотожно рівний вираз}
\matchItem{А}{$a + \sqrt{2}$}
\matchItem{Б}{$1$}
\matchItem{В}{$a + 3$}
\matchItem{Г}{$a + 3b$}
\matchItem{Д}{$a - \sqrt{2}$}
}{
\matchingGrid
}
\par\penalty-20\vspace{0.25cm}
\end{samepage}

\begin{samepage}
% Завдання 1 (на відповідність)
\zadtask{До кожного виразу \mbox{(1--3)} доберіть тотожно рівний йому вираз \mbox{(А--Д)}, якщо $a \neq 8$. \nmtyear{2023}}

\nmtnobreak
\nopagebreak\nbvspace{0.3cm}
\matchingLayout{
\matchHead{Вираз}
\matchItem{1}{$\dfrac{64 - a^2}{8 - a}$}
\matchItem{2}{$2^{3a}$}
\matchItem{3}{$\log_{\sqrt[8]{2}} 2^a$}
}{
\matchHead{Тотожно рівний вираз}
\matchItem{А}{$8^a$}
\matchItem{Б}{$8 + a$}
\matchItem{В}{$8a$}
\matchItem{Г}{$8 - a$}
\matchItem{Д}{$\dfrac{a}{8}$}
}{
\matchingGrid
}
\par\penalty-20\vspace{0.25cm}
\end{samepage}

\begin{samepage}
% Завдання 3 (на відповідність)
\zadtask{До кожного виразу \mbox{(1--3)} доберіть тотожно рівний йому вираз \mbox{(А--Д)}, якщо $a > 0$. \nmtyear{2023}}

\nmtnobreak
\nopagebreak\nbvspace{0.3cm}
\matchingLayout{
\matchHead{Вираз}
\matchItem{1}{$\sqrt{4a}$}
\matchItem{2}{$2^{\log_4 a}$}
\matchItem{3}{$\left(\dfrac{2}{a}\right)^{-1}$}
}{
\matchHead{Тотожно рівний вираз}
\matchItem{А}{$-\dfrac{2}{a}$}
\matchItem{Б}{$2a$}
\matchItem{В}{$2\sqrt{a}$}
\matchItem{Г}{$\sqrt{a}$}
\matchItem{Д}{$\dfrac{a}{2}$}
}{
\matchingGrid
}
\par\penalty-20\vspace{0.25cm}
\end{samepage}

\begin{samepage}
% НМТ 2026, сесія 3.06, завдання №16 --- основна тема: Квадратні корені. Корені вищих степенів. Раціональний степінь
\noindent\zadnum Узгодьте вираз \mbox{(1--3)} з тотожно рівним йому виразом \mbox{(А--Д)}, якщо $a \neq 0$. \nmtyear{2026}

\nmtnobreak
\nopagebreak\nbvspace{0.3cm}
\noindent
\matchingLayout{
\matchHead{Вираз}
\matchItem{1}{$\sqrt[3]{\dfrac{27}{a^3}}$}
\matchItem{2}{$a^{-2}\cdot a^3 \cdot \log_{27} 3$}
\matchItem{3}{$\dfrac{(a-3)^2 - a^2 - 9}{2}$}
}{
\matchHead{Тотожно рівний вираз}
\matchItem{А}{$-3a$}
\matchItem{Б}{$\dfrac{3}{a}$}
\matchItem{В}{$0$}
\matchItem{Г}{$\dfrac{a}{3}$}
\matchItem{Д}{$3a$}
}{
\matchingGrid
}
% Відповідь: 1~--~Б, 2~--~Г, 3~--~А
\ifshowsolutions\solution{\textbf{1}: $\sqrt[3]{\dfrac{27}{a^3}} = \dfrac{3}{a}$~--- Б. \textbf{2}: $a^{-2}\cdot a^3\cdot\log_{27}3 = a\cdot\dfrac13 = \dfrac{a}{3}$~--- Г. \textbf{3}: $\dfrac{(a-3)^2 - a^2 - 9}{2} = \dfrac{-6a}{2} = -3a$~--- А. Відповідь: \textbf{1~--~Б, 2~--~Г, 3~--~А}.}\fi
\par\penalty-20
\end{samepage}

\begin{samepage}
% НМТ 2026, сесія 12.06, завдання №16 --- також у темі: Многочлени. Формули скороченого множення, Логарифмічні вирази
\zadtask{Узгодьте вираз \mbox{(1--3)} з тотожно рівним йому виразом \mbox{(А--Д)}. \nmtyear{2026}\par
\nopagebreak\nbvspace{0.3cm}
\noindent
\matchingLayout{
\matchHead{Числовий вираз}
\matchItem{1}{$\dfrac{1}{n} : \dfrac{1}{2}$}
\matchItem{2}{$10n^2 - 2n(5n - 1)$}
\matchItem{3}{$\log_{25} 5^n$}
}{
\matchHead{Тотожно рівний вираз}
\matchItem{А}{$-2n$}
\matchItem{Б}{$\dfrac{2}{n}$}
\matchItem{В}{$\dfrac{n}{2}$}
\matchItem{Г}{$2n$}
\matchItem{Д}{$2n^2$}
}{
\matchingGrid
}
}
\par\penalty-20
\end{samepage}

\begin{samepage}
\needspace{0.45\textheight}\typeTitle{Числовий вираз, рівний даному: ірраціональність у знаменнику, модуль}
% Завдання 10 (на відповідність)
\zadtask{До кожного виразу \mbox{(1--3)} доберіть тотожно рівний йому вираз \mbox{(А--Д)}. \nmtyear{2023}}

\nmtnobreak
\nopagebreak\nbvspace{0.3cm}
\noindent
\matchingLayout{
\matchHead{Вираз}
\matchItem{1}{$|1 - \sqrt{5}| - \sqrt{5} + 1$}
\matchItem{2}{$\dfrac{2\sqrt{5} - 10}{\sqrt{5}}$}

\nmtnobreak
\matchItem{3}{$\log_{\sqrt{5}} 5$}
}{
\matchHead{Тотожно рівний вираз}
\matchItem{А}{$\sqrt{5}$}
\matchItem{Б}{$0$}
\matchItem{В}{$2 - 2\sqrt{5}$}
\matchItem{Г}{$2$}
\matchItem{Д}{$-8$}
}{
\matchingGrid
}
\par\penalty-20\vspace{0.25cm}
\end{samepage}

\begin{samepage}
% Завдання 20 (на відповідність)
\zadtask{До кожного виразу \mbox{(1--3)} доберіть тотожно рівний йому вираз \mbox{(А--Д)}. \nmtyear{2023}}

\nmtnobreak
\nopagebreak\nbvspace{0.3cm}
\matchingLayout{
\matchHead{Вираз}
\matchItem{1}{$\dfrac{1}{\sqrt{10} - 3}$}

\nmtnobreak
\matchItem{2}{$|3 - \sqrt{10}|$}
\matchItem{3}{$\log_5 125$}
}{
\matchHead{Тотожно рівний вираз}
\matchItem{А}{$\sqrt{10} - 3$}
\matchItem{Б}{$3 - \sqrt{10}$}
\matchItem{В}{$\sqrt{10} + 3$}
\matchItem{Г}{$3$}
\matchItem{Д}{$25$}
}{
\matchingGrid
}
\par\penalty-20\vspace{0.25cm}
\end{samepage}

\begin{samepage}
\needspace{0.45\textheight}\typeTitle{Початок речення: вираз $a$ через $n$}
\noindent\zadnum До кожного початку речення (1–3) доберіть його закінчення (А–Д) так, щоб утворилося правильне твердження, якщо $n > 0$. \nmtyear{2023}\par
\nopagebreak\nbvspace{0.3cm}

\nmtnobreak
\noindent
\matchingLayout{
\matchHead{Початок речення}
    \matchItem{1}{Якщо $\displaystyle\frac{n}{a} = 3$, то}
\matchItem{2}{Якщо $1 + \log_3 n = \log_3 a$, то}
\matchItem{3}{Якщо $3^n \cdot 3 = 3^a$, то}
}{
\matchHead{Закінчення речення}
    \matchItem{А}{$a = 3n$.}
\matchItem{Б}{$a = n + 1$.}
\matchItem{В}{$a = n + 3$.}
\matchItem{Г}{$a = \displaystyle\frac{3}{n}$.}
\matchItem{Д}{$a = \displaystyle\frac{n}{3}$.}
}{
\matchingGrid
}
\par\penalty-20\vspace{0.25cm}
\end{samepage}

\begin{samepage}
% Завдання 3 (на відповідність)
\zadtask{До кожного початку речення \mbox{(1--3)} доберіть його закінчення \mbox{(А--Д)} так, щоб утворилося правильне твердження, якщо $n > 0$. \nmtyear{2023}}

\nmtnobreak
\nopagebreak\nbvspace{0.3cm}
\noindent
\matchingLayout{
\matchHead{Початок речення}
\matchItem{1}{Якщо $|-n| = a$, то}
\matchItem{2}{Якщо $(n - 2)(n + 1) - n^2 = a$, то}
\matchItem{3}{Якщо $2^{-\lg n} = 2^{\lg a}$, то}
}{
\matchHead{Закінчення речення}
\matchItem{А}{$a = -n - 2$.}
\matchItem{Б}{$a = \dfrac{1}{n}$.}
\matchItem{В}{$a = n$.}
\matchItem{Г}{$a = -2$.}
\matchItem{Д}{$a = -n$.}
}{
\matchingGrid
}
\par\penalty-20\vspace{0.25cm}
\end{samepage}

\begin{samepage}
\noindent\zadnum До кожного початку речення \mbox{(1--3)} доберіть його закінчення \mbox{(А--Д)} так, щоб утворилося правильне твердження, якщо $m$ і $n$ -- натуральне число, $n > 1$, $m > 1$. \nmtyear{2023}\par
\nopagebreak\nbvspace{0.3cm}

\nmtnobreak
\matchingLayout{
\matchHead{Початок речення}

\nmtnobreak
\matchItem{1}{Якщо $n^{\log_n m} = a$, то}
\matchItem{2}{Якщо $\displaystyle\frac{mn^2 - nm^2}{n - m} = a$, то}
\matchItem{3}{Якщо $n\cos^2 m + n\sin^2 m = a$, то}
}{
\matchHead{Закінчення речення}

\nmtnobreak
\matchItem{А}{$a = m + n$.}
\matchItem{Б}{$a = m$.}
\matchItem{В}{$a = m - n$.}
\matchItem{Г}{$a = mn$.}
\matchItem{Д}{$a = n$.}
}{
\matchingGrid
}
\par\penalty-20\vspace{0.25cm}
\end{samepage}

\begin{samepage}
% Завдання 7 (на відповідність)
\zadtask{До кожного початку речення \mbox{(1--3)} доберіть його закінчення \mbox{(А--Д)} так, щоб утворилося правильне твердження, якщо $m$ і $n$ --- натуральні числа, $n > 1$, $m > 1$. \nmtyear{2023}}

\nmtnobreak
\nopagebreak\nbvspace{0.3cm}
\matchingLayout{
\matchHead{Початок речення}
\matchItem{1}{Якщо $n \sin m\pi = a$, то}
\matchItem{2}{Якщо $\dfrac{2^m}{2^n} = 2^a$, то}
\matchItem{3}{Якщо $\sqrt[n]{\sqrt[m]{2}} = \sqrt[a]{2}$, то}
}{
\matchHead{Закінчення речення}
\matchItem{А}{$a = mn$.}
\matchItem{Б}{$a = 0$.}
\matchItem{В}{$a = m - n$.}
\matchItem{Г}{$a = n$.}
\matchItem{Д}{$a = \dfrac{m}{n}$.}
}{
\matchingGrid
}
\par\penalty-20\vspace{0.25cm}
\end{samepage}

\begin{samepage}
\noindent\zadnum До кожного початку речення (1–3) доберіть його закінчення (А–Д) так, щоб утворилося правильне твердження. \nmtyear{2023}\par
\nopagebreak\nbvspace{0.3cm}

\nmtnobreak
\noindent
\matchingLayout{
\matchHead{Початок речення}
    \matchItem{1}{Якщо $3^m \cdot 9^n = 3^a$, то}
\matchItem{2}{Якщо $(m + n)^2 - m^2 - n^2 = a$, то}
\matchItem{3}{Якщо $\log_3 27^{mn} = a$, то}
}{
\matchHead{Закінчення речення}
    \matchItem{А}{$a = 0$.}
\matchItem{Б}{$a = 27^{mn}$.}
\matchItem{В}{$a = 3mn$.}
\matchItem{Г}{$a = 2mn$.}
\matchItem{Д}{$a = m + 2n$.}
}{
\matchingGrid
}
\par\penalty-20\vspace{0.25cm}
\end{samepage}

\begin{samepage}
% Завдання 4 (на відповідність)
\zadtask{До кожного початку речення \mbox{(1--3)} доберіть його закінчення \mbox{(А--Д)} так, щоб утворилося правильне твердження, якщо $n$ --- натуральне число, $n > 1$. \nmtyear{2023}}

\nmtnobreak
\nopagebreak\nbvspace{0.3cm}
\matchingLayout{
\matchHead{Початок речення}
\matchItem{1}{Якщо $n \cos 8\pi = a$, то}
\matchItem{2}{Якщо $\log_2 8 + \log_2 n = \log_2 a$, то}
\matchItem{3}{Якщо $\sqrt{\sqrt[n]{8}} = \sqrt[a]{8}$, то}
}{
\matchHead{Закінчення речення}
\matchItem{А}{$a = 2n$.}
\matchItem{Б}{$a = 8n$.}
\matchItem{В}{$a = 8 + n$.}
\matchItem{Г}{$a = n$.}
\matchItem{Д}{$a = 3n$.}
}{
\matchingGrid
}
\par\penalty-20\vspace{0.25cm}
\end{samepage}

\begin{samepage}
% === ЗАВДАННЯ 6 (Відповідність) ===
\noindent\zadnum До кожного початку речення \mbox{(1--3)} доберіть його закінчення \mbox{(А--Д)} так, щоб утворилося правильне твердження, якщо $n$ -- натуральне число, $n>1$. \nmtyear{2023}

\nmtnobreak
\nopagebreak\nbvspace{0.3cm}

\nmtnobreak
\noindent
\matchingLayout{
\matchHead{Початок речення}
    \matchItem{1}{\trigStyle{Якщо $n \cos 8\pi = a$, то}}
\matchItem{2}{Якщо $\log_2 8 + \log_2 n = \log_2 a$, то}
\matchItem{3}{Якщо $\sqrt[n]{\sqrt[3]{8}} = \sqrt[a]{8}$, то}
}{
\matchHead{Закінчення речення}
    \matchItem{А}{$a = 2n$.}
\matchItem{Б}{$a = 8n$.}
\matchItem{В}{$a = 8 + n$.}
\matchItem{Г}{$a = n$.}
\matchItem{Д}{$a = 3n$.}
}{
\matchingGrid
}
\par\penalty-20\vspace{0.25cm}
\end{samepage}

\begin{samepage}
% Завдання 38 (на відповідність)
\zadtask{Доберіть до кожного початку речення \mbox{(1--3)} його закінчення \mbox{(А--Д)} так, щоб утворилося правильне твердження, якщо $n \neq 0$. \nmtyear{2025}}

\nmtnobreak
\nopagebreak\nbvspace{0.3cm}
\noindent
\matchingLayout{
\matchHead{Початок речення}
\matchItem{1}{Якщо $(m + 1)^2 - 1 = n$, то}
\matchItem{2}{Якщо $\sqrt{2^m} = 2^n$, то}
\matchItem{3}{Якщо $\log_{2^n} 4^m = 1$, то}
}{
\matchHead{Закінчення речення}
\matchItem{А}{$n = 2m$.}
\matchItem{Б}{$n = m^2$.}
\matchItem{В}{$n = m^2 + 2m$.}
\matchItem{Г}{$n = \sqrt{m}$.}
\matchItem{Д}{$n = \dfrac{m}{2}$.}
}{
\matchingGrid
}
\par\penalty-20\vspace{0.25cm}
\end{samepage}

\begin{samepage}
% Завдання 51 (на відповідність)
\zadtask{Доберіть до кожного початку речення \mbox{(1--3)} його закінчення \mbox{(А--Д)} так, щоб утворилося правильне твердження, якщо $m > 0$. \nmtyear{2025}}

\nmtnobreak
\nopagebreak\nbvspace{0.3cm}
\matchingLayout{
\matchHead{Початок речення}
\matchItem{1}{Якщо $m \cos x \cdot \mathrm{tg}\, x = 2n \sin x$, то}
\matchItem{2}{Якщо $\sqrt{2m^2} = n$, то}
\matchItem{3}{Якщо $4^{\log_2 m} = n$, то}
}{
\matchHead{Закінчення речення}
\matchItem{А}{$n = 2m$.}
\matchItem{Б}{$n = m^2$.}
\matchItem{В}{$n = 2^m$.}
\matchItem{Г}{$n = m\sqrt{2}$.}
\matchItem{Д}{$n = \dfrac{m}{2}$.}
}{
\matchingGrid
}
\par\penalty-20\vspace{0.25cm}
\end{samepage}

\begin{samepage}
\needspace{0.45\textheight}\typeTitle{Значення змінної, за якого вираз дорівнює $1$}
% Завдання 7 (на відповідність)
\zadtask{Узгодьте вираз \mbox{(1--3)} із значенням $m$ \mbox{(А--Д)}, за якого значення цього виразу дорівнює $1$. \nmtyear{2025}}

\nmtnobreak
\nopagebreak\nbvspace{0.3cm}
\matchingLayout{
\matchHead{Вираз}
\matchItem{1}{$\dfrac{m}{4}$}
\matchItem{2}{$4^6 : 4^{-m}$}
\matchItem{3}{$\log_{16} 2 + \log_{16} m$}
}{
\matchHead{Значення $m$}
\matchItem{А}{$\dfrac{1}{4}$}
\matchItem{Б}{$6$}
\matchItem{В}{$4$}
\matchItem{Г}{$-6$}
\matchItem{Д}{$8$}
}{
\matchingGrid
}
\par\penalty-20\vspace{0.25cm}
\end{samepage}

\clearpage
\sectionTitle{Відповіді}

\noindent{\small\color{gray!80!black}Для кожного завдання -- літери відповідей для пунктів 1, 2, 3. Відповіді обчислено й звірено з офіційними ключами там, де вони є.}\par\vspace{0.3cm}

\noindent\textbf{Частина 1}\par\vspace{0.1cm}

\begin{multicols}{5}\noindent
\mbox{1~--- ВБА}\par
\mbox{2~--- ГВБ}\par
\mbox{3~--- АДГ}\par
\mbox{4~--- АДГ}\par
\mbox{5~--- ДВГ}
\end{multicols}\par\vspace{0.2cm}

\noindent\textbf{Частина 2}\par\vspace{0.1cm}

\begin{multicols}{5}\noindent
\mbox{6~--- ВДБ}\par
\mbox{7~--- ДГВ}\par
\mbox{8~--- АГБ}\par
\mbox{9~--- ГАВ}\par
\mbox{10~--- ДБА}\par
\mbox{11~--- ДБГ}\par
\mbox{12~--- ВГБ}\par
\mbox{13~--- ГБВ}\par
\mbox{14~--- АДВ}\par
\mbox{15~--- ВБА}\par
\mbox{16~--- ГАВ}\par
\mbox{17~--- АГВ}\par
\mbox{18~--- ГАВ}\par
\mbox{19~--- АБВ}\par
\mbox{20~--- ГАВ}\par
\mbox{21~--- ВБД}\par
\mbox{22~--- АДБ}\par
\mbox{23~--- АВД}\par
\mbox{24~--- ГВБ}\par
\mbox{25~--- ВГБ}\par
\mbox{26~--- БАВ}\par
\mbox{27~--- ДГВ}\par
\mbox{28~--- ДВБ}\par
\mbox{29~--- ДВБ}\par
\mbox{30~--- ДБВ}\par
\mbox{31~--- БДВ}\par
\mbox{32~--- ГВА}\par
\mbox{33~--- АГВ}\par
\mbox{34~--- ВБД}\par
\mbox{35~--- ГВБ}\par
\mbox{36~--- АБД}\par
\mbox{37~--- БДГ}\par
\mbox{38~--- ГАВ}\par
\mbox{39~--- ВГБ}\par
\mbox{40~--- ВБД}\par
\mbox{41~--- ДГБ}\par
\mbox{42~--- ДГА}\par
\mbox{43~--- ГДВ}\par
\mbox{44~--- ДГВ}\par
\mbox{45~--- ВДБ}\par
\mbox{46~--- ДВБ}\par
\mbox{47~--- ГДА}\par
\mbox{48~--- ВГБ}\par
\mbox{49~--- ВБА}\par
\mbox{50~--- АВГ}\par
\mbox{51~--- ДГБ}\par
\mbox{52~--- АГВ}\par
\mbox{53~--- БАГ}\par
\mbox{54~--- АБГ}\par
\mbox{55~--- БВА}\par
\mbox{56~--- АБГ}\par
\mbox{57~--- ДАВ}\par
\mbox{58~--- БАВ}\par
\mbox{59~--- ВГД}\par
\mbox{60~--- БГА}\par
\mbox{61~--- БГВ}\par
\mbox{62~--- БВГ}\par
\mbox{63~--- ВАГ}\par
\mbox{64~--- ДАБ}\par
\mbox{65~--- ВАБ}\par
\mbox{66~--- БГД}\par
\mbox{67~--- БВА}\par
\mbox{68~--- ДГВ}\par
\mbox{69~--- ГБА}\par
\mbox{70~--- ГБД}\par
\mbox{71~--- ВДА}\par
\mbox{72~--- ДГБ}\par
\mbox{73~--- ВГД}
\end{multicols}\par\vspace{0.2cm}

\end{document}
